\pdfoutput=1
% ^ MUST stay within the first five lines, and comments count toward that budget. arXiv's TeX
% pipeline decides between latex (DVI) and pdflatex by inspection, and a DVI run cannot
% \includegraphics a PDF figure; without this declaration a submission carrying PDF figures can be
% announced with its figures silently missing.
\documentclass[11pt]{article}

% ---------------------------------------------------------------------------
% Paper I of three: The Conditional Lever
% Build: pdflatex main && bibtex main && pdflatex main x2
% ---------------------------------------------------------------------------
% Shared preamble for the three papers of this series. Each paper's main.tex declares
% \pdfoutput and \documentclass itself -- \pdfoutput must sit within the first five lines of the
% file arXiv reads, and an \input cannot satisfy that -- then inputs this file for everything else.
%
% Kept as one file rather than triplicated because every setting below was arrived at by fixing a
% specific rendering defect, and three copies would drift apart at the first fix.

\usepackage[utf8]{inputenc}
\usepackage[T1]{fontenc}
\usepackage{amsmath,amssymb}
\usepackage{booktabs}
\usepackage{graphicx}
\usepackage{hyperref}
% Default hyperref draws a coloured rule around every link, so every \S and \ref printed as a red
% box. Invisible in the source and unmissable on paper; found by rasterising a page rather than by
% reading the log, which reports nothing about it.
\hypersetup{hidelinks}
\usepackage{natbib}
\usepackage[margin=1in]{geometry}
\usepackage{xcolor}
% Long \texttt{} identifiers and inline result vectors give TeX nothing to break on, which
% cost eleven overfull boxes running up to 1.7in into the margin. This lets it stretch
% interword space as a last resort instead; content is unchanged.
\setlength{\emergencystretch}{3em}
% Float placement. With many tables queued, LaTeX's defaults (a top float may occupy at most 0.7 of
% the text height, and 0.2 of every float-bearing page must be text) deferred two full-width figures
% to the end of the document --- twenty-eight pages from the section that references them. Relaxing
% the fractions lets a tall figure sit at the top of its own page near its reference instead of
% falling off the back of the float queue.
\renewcommand{\topfraction}{0.92}
\renewcommand{\bottomfraction}{0.75}
\renewcommand{\textfraction}{0.08}
\renewcommand{\floatpagefraction}{0.70}
\setcounter{topnumber}{3}
\setcounter{totalnumber}{4}
\newcommand{\dd}{\ensuremath{\Delta}}


\title{Whether Optimizing the Loop Beats Optimizing the Prompt\\
Depends on One Number\\
\large $p_0$, the base rate of already-correct first drafts,\\
and a $0.102$ swing in the contrast across it}

\author{Vladimir Rotariu \\ Independent researcher \\
  \texttt{github.com/loopgepa/whystuck}}
\date{July 2026}

\begin{document}
\maketitle

\begin{abstract}
Reflective optimizers such as GEPA \citep{agrawal2025gepa} can evolve several text components of a
compound LLM system, so a natural hypothesis is that optimizing an entire
generate--critique--refine loop beats optimizing its generation prompt alone. On a corrected,
control-heavy harness the answer turns on one quantity: $p_0$, the base rate of first drafts
already satisfying every requirement. At $p_0 = 0.067$ the search converges on a verifier that
never accepts anything, on $11$ of $12$ evolved genes across four training-split sizes, and
evolving it buys nothing measurable; at $p_0 = 0.267$, changing only which constraints the
generator may use, evolving the verifier alongside the generator \emph{gains} $0.060$ held-out,
unanimous across seeds and $p<0.10$ over five independent draws. Neither endpoint carries the
result on its own: what does is the $0.102$ swing between them, about $2.5\sigma$ of measured
noise, along an axis the mechanism named before the runs were bought. Unconditional refinement is
nearly free on a training signal where almost nothing is already correct and costly where much of
it is, which makes the degenerate verifier optimal at low $p_0$ and screenable, before spending, as
a \texttt{PASS} rate of zero. Two pre-registered sweeps close the alternative readings: no
split-size trend to overfit along from $4$ to $24$ training tasks, and no trend over a measured
six-fold range of search depth, where the deepest cell is the flattest while \emph{acceptance}
rises with budget. Reaching these measurements meant
discarding fifteen earlier contrast runs,
whose critique gene the optimizer never touched (companion Paper~III).
We also show $\dd$ is a whole-system contrast rather than a treatment effect: decomposed, $60\%$ of
the low-$p_0$ deficit is the gene both arms share while all of the high-$p_0$ gain is the
treatment. Every published claim this paper withdraws is listed in
Appendix~\ref{sec:corrections}.
\end{abstract}

% The companion-paper note. Every paper in the series carries it so a reader who arrives at one
% knows what the other two contain and, more importantly, which claims are NOT in the paper they
% are holding. Three papers were carved out of one 46-page draft; the carve is documented in
% docs/papers/README.md, and the split's completeness is machine-checked by
% scripts/verify_split_coverage.py against the pre-split source.

\paragraph{This paper is one of three.}
One study produced three separable results, and they are reported separately because they answer
different questions, rest on different experiments, and are useful to different readers. All three
share one harness, one artifact release and one corrections register.
\begin{description}\itemsep0pt
  \item[I. The conditional lever.] Whether optimizing a self-refine loop's verifier beats
    optimizing its generation prompt, and the quantity that decides it: $p_0$, the base rate of
    first drafts already satisfying every requirement.
  \item[II. What a checker cannot buy is the repair.] The loop's own ceiling, with no optimizer
    anywhere in it: how good the intrinsic critique is as a detector, what happens when the
    detector is made perfect, and which fix-side levers actually pay.
  \item[III. Your loop optimizer is not optimizing your loop.] A silent failure mode in
    componentized prompt optimization: pointing an optimizer at a component is not the same as
    optimizing it, the difference does not appear in held-out scores, and one integer per component
    catches it.
\end{description}
Shared artifact, including every cassette, pre-registration and the corrections register:
\texttt{github.com/loopgepa/whystuck}.


\section{Introduction}
\label{sec:intro}

Automatic prompt optimization (APO) has matured from single-instruction search
\citep{madaan2023selfrefine} to optimizers that treat an entire LLM program as
the object of optimization. GEPA \citep{agrawal2025gepa}, the current
state-of-the-art reflective optimizer and an ICLR-2026 oral, evolves any set of
text components of a compound system --- prompts, rubrics, tool descriptions ---
by reflecting on execution traces in natural language and maintaining a Pareto
frontier of candidates. DSPy \citep{khattab2023dspy} similarly compiles and
optimizes multi-module pipelines end-to-end.

This capability raises a concrete engineering question that, to our knowledge,
has not been isolated experimentally for a self-refine loop's own verifier. The
adjacent question for multi-\emph{agent} systems has since been benchmarked:
\citet{bai2026masprompt} find that GEPA improves single-agent performance across all five of
their tasks, by $+4.2$ points on average, while its multi-agent extension ranges from $-16.0$ to
$+24.0$ points depending on task and workflow topology --- more optimizable surface is not reliably
better. Our question is narrower and about a different object: \emph{given a self-refine loop, does
optimizing the loop's verification/critique component (and beyond) yield better
held-out performance than optimizing only its generation prompt, under a matched
budget?} We call the former \textbf{loop-level} and the latter
\textbf{prompt-only} optimization. The hypothesis is intuitive --- more
optimizable surface, more headroom --- and if true it would justify the extra
cost of evolving critique rubrics and plan modules.

The answer we arrive at is neither yes nor no but \emph{conditional}, and the condition is a
property of the task distribution rather than of the optimizer. Write $p_0$ for the base rate of
first drafts that already satisfy every constraint. A verifier that never accepts sends every draft
to a rewrite, which costs only on drafts that needed none, and the weight of that stratum is
exactly $p_0$. So the laziest available verifier is nearly free where $p_0$ is small and expensive
where it is large, and a search that optimizes a training signal will take it in the first case and
not the second. That is a quantitative commitment made before the runs, and it is what turns a null
into a law.

We answer the question with a pre-registered, control-heavy study. Our
contributions:
\begin{enumerate}
  \item \textbf{The measurement, and the condition it depends on} (\S\ref{sec:results}): under
    a matched budget on a corrected harness, loop-level optimization does not beat prompt-only
    optimization at low $p_0$, with the treatment
    verified to have reached the gene on every seed. Holding harness, optimizer, budget and
    seeds fixed and changing only which constraints the task generator may use raises $p_0$ from
    $0.067$ to $0.267$, and the comparison turns to $+0.060$ with all three seeds agreeing
    (\S\ref{sec:p0}). That is the study's largest loop-level effect and the one that replicates.
    A first draft of this paper also reported the low-$p_0$ cell as a \emph{cost} of $0.042$;
    a pre-registered split-size sweep and an offline estimate of the study's own measurement
    noise together withdraw that, and \S\ref{sec:splitresult} reports the withdrawal rather than
    burying it. What the condition governs, robustly, is what the optimizer does to the
    verifier; what it costs is smaller than we could measure.
  \item \textbf{The swing, not the endpoint, is the finding} (\S\ref{sec:p0}). Neither cell is
    individually significant at $n=3$ seeds, where an exact sign test cannot fall below
    $p = 0.25$. What noise does not account for is the difference between the levels,
    $0.102$, which is about $2.5\sigma$ once the denominator is allowed to carry $\sigma$ twice
    --- and which the mechanism predicted in advance.
  \item \textbf{The mechanism} (\S\ref{sec:analysis}): the optimizer converges on a verifier
    that never accepts. Scored as a detector, the evolved gene emits \texttt{PASS} zero times
    on every seed at low $p_0$ and loses the one true positive the seed rubric had.
    Unconditional refinement is near-free on a training signal where most first drafts fail,
    so this is Goodhart's law operating inside the loop --- and it is screenable, before
    spending, as a \texttt{PASS} rate of zero. As $p_0$ rises the same search learns to accept,
    and does so at \emph{precision~1}: it never emits a false \texttt{PASS} at any level, while
    the seed rubric it is compared against emits six.
  \item \textbf{How much of a componentized $\dd$ is the component} (\S\ref{sec:decomp}):
    because both arms evolve the shared \texttt{generate} gene, $\dd(D-B)$ is a whole-system
    contrast and not a measurement of the treatment. Splitting it into a draft term and a loop
    term --- an identity, computable offline from the released artifact --- shows the $-0.042$
    is $60\%$ shared gene while the $+0.060$ is $100\%$ treatment. It also bounds the
    outcome-side half of the mechanism: the stratum the always-\texttt{FIX} account blames
    contributes $+0.002$, not the deficit.
  \item \textbf{Two pre-registered sweeps that close the alternative readings}
    (\S\ref{sec:splitresult}, \S\ref{sec:budgetresult}). Small-split overfitting was our own
    first explanation and it does not survive: across $n_{\mathrm{train}} \in \{4,8,16,24\}$ at
    constant epochs there is no split-size trend, three of four cells are positive, and the one
    negative cell is the run we led with. Under-search does not survive either: over a
    \emph{measured} six-fold depth range $\dd$ does not order along depth and the deepest cell is
    the flattest, while acceptance rises with budget --- which separates \emph{the optimizer
    failed to search} from \emph{the optimizer searched and found nothing that transfers}.
  \item \textbf{A positive control} (\S\ref{sec:posctrl}), because every headline here is a
    null: on a landscape where the critique gene is the only path to a passing answer, the
    same harness and optimizer return $\dd(D-B) = +1.000$.
\end{enumerate}

We report every result as-is, including a run in which optimization \emph{degraded} held-out
performance, and including the audit that invalidated our own first headline. That audit is the
subject of Paper~III and is the reason the measurements below exist at all: our first fifteen
contrast runs measured $\dd$ from an arm that never received its treatment. The point is not
that GEPA is weak --- it is a strong optimizer --- but that \emph{pointing an optimizer at a
component is not the same as optimizing it}.

\section{Background and Related Work}
\label{sec:related}

\paragraph{Loop engineering, and which loop this paper is about.}
During June 2026 industry practice converged on a name for designing the machinery
\emph{around} a model rather than its prompt. Steinberger compressed it into a slogan ---
stop prompting coding agents, design the loops that prompt them
\citep{steinberger2026loops} --- Osmani's essay named the practice \emph{loop engineering}
\citep{osmani2026loop}, and Swyx framed the general form as \emph{loopcraft}, ``the art of
stacking loops'' \citep{swyx2026loopcraft}. Of the four loops that decomposition names
\citep{langchain2026loopcraft} --- an \textbf{agent} loop calling tools until done, a
\textbf{verification} loop scoring output against a rubric and retrying with feedback, an
\textbf{event-driven} loop triggering runs from outside, and a \textbf{hill-climbing} loop that
feeds execution traces to an analysis step which improves the harness --- two are ours and two are
not. The verification loop is our object of study: the \texttt{ensure\_correct} gene \emph{is} the
verifier, the PASS gate plus \texttt{loop\_max\_retries} \emph{is} the stop rule. The hill-climbing
loop, called ``arguably most important'' \citep{langchain2026loopcraft}, is our intervention: GEPA
reflects on recorded traces and rewrites the loop's own text components, so arm~D is a
hill-climbing loop pointed at a verification loop. The agent and event-driven loops are out of
scope, since our worker is a fixed generate--critique--refine sequence with no tool-calling agent
and is invoked once per benchmark task, with no triggers, schedules or cross-task memory.

So the paper's question, in this vocabulary, is: \emph{can a hill-climbing loop improve a
verification loop by rewriting its verifier?} ``Loop-level'' throughout means the verifier
and its stop rule --- not the full stack --- and we vary topology only through an optional
\texttt{plan} module and an optional repair \texttt{actuator}. Chronology: our runs were
pre-registered and executed around and before these posts appeared, so the correspondence is
convergent rather than causal --- the practitioner literature supplied the vocabulary and the
stakes, not the hypotheses.

\paragraph{Reflective and evolutionary prompt optimization.}
GEPA \citep{agrawal2025gepa} reports beating GRPO by $6\%$ on average and up to $20\%$ while
using up to $35\times$ fewer rollouts, and exceeding MIPROv2 by over $10\%$. Crucially for us, GEPA
optimizes \emph{one or more} components and is known to plateau. The sharpest published
instance is independent of us: optimizing a prompt with GEPA across six tasks,
\citet{pappu2026multiagent} report that ``on GPQA Diamond and MATH-500, GEPA returned the
original prompt essentially unchanged.'' OPRO \citep{yang2023opro} and EvoPrompt
\citep{guo2024evoprompt} likewise use an
LLM as the optimizer, and OPRO's authors quantify the over-fitting risk directly: their
``training accuracies are often 5\%--20\% higher than our test accuracies''. Our null is a
member of this ``plateau/degrade'' family, isolated under controlled conditions --- and our
contribution against it is the \emph{condition}: we exhibit a task-distribution parameter that
moves the same optimizer, harness and budget from no measurable effect to the study's largest
gain.

\paragraph{Self-refine and the limits of self-correction.}
Self-Refine \citep{madaan2023selfrefine} popularized generate--critique--refine
with a single model, reporting ${\sim}20\%$ absolute improvement on average across seven tasks.
Huang et al.\ \citep{huang2023cannot} show the sharp boundary: \emph{intrinsic}
self-correction (no external feedback) fails to improve and often degrades, and prior positive
results relied on \emph{oracle} stopping signals that ``vanish when oracle labels are
not available.'' A parallel line shows LLM
self-verifiers are unreliable detectors --- high false-positive rates
\citep{valmeekam2023selfcritique} and a self-enhancement bias, with cross-family
verification outperforming self-verification. The boundary is not undisputed:
\citet{liu2024intrinsic} report intrinsic self-correction working under narrow conditions
(zero temperature, prompts free of leading bias), which is a condition our worker does not
satisfy --- it samples at the provider's default temperature (\S\ref{sec:setup}), and that
choice is what makes the noise floor of \S\ref{sec:threats} non-zero in the first place.
These results predict our loop-level null, and Paper~II measures the loop's own ceiling
independently of any optimizer: the intrinsic critique's discrimination does not improve on
base-rate-independent statistics as the tasks get easier, even as its $\mathrm{F}_1$ nearly triples,
and a \emph{perfect} inference-time checker does not beat it.

\paragraph{When does optimization help?}
Recent work explicitly frames the boundary question for multi-agent systems and
cost-aware settings. We adopt the same posture for the narrower loop-level vs.\
prompt-only question and contribute a pre-registered answer with controls.

\section{Method}
\label{sec:method}

% The harness description common to all three papers of this series. Copied into each rather than
% cited across them, because a reader of any one paper needs it to read that paper's numbers.
%
% Kept as a single input, verbatim, for a reason specific to this study: the claim checker
% (scripts/verify_paper_claims.py) matches several of these sentences as literal text, spaces and
% line breaks included, so three hand-maintained copies would drift and silently stop being
% checked. Edit here, and never reflow a line that contains a number.

\paragraph{The loop.}
The worker is a real generate--critique--refine loop. A \texttt{generate} gene, which the prose also calls the
generation gene, drafts an answer; a \texttt{ensure\_correct} gene, which it calls the
verify, verifier or critique gene interchangeably, drives an LLM
self-critique that returns \texttt{PASS} or a list of fixes and gates a bounded
number of refine rounds; the final answer is scored by a programmatic metric. Two further
genes are optional and named where they are used: a \texttt{plan} module that drafts before
generating, and a repair \texttt{actuator} that turns a critique into the instructions the
refine step is handed.
Critically, the critique is \emph{intrinsic} wherever the loop is left as built: the
programmatic scorer is never consulted inside the loop (that would be oracle contamination), so
the loop cannot game its own metric. Paper~II suspends this deliberately in its oracle-critique
controls, which hand the in-loop gate the true metric and are named as such wherever they are
reported; nothing else in the series does.
Each eval costs \emph{at most} $1 + 2R$ LLM calls for a refine
budget of $R$ (or $2 + 2R$ with the optional plan module), and fewer whenever a critique passes
before the budget is spent: the loop stops on a \texttt{PASS} without issuing the refine call
that would have followed it, so the bound is reached only by an eval that never passes.

\paragraph{Tasks.} A composable IFEval-style \citep{zhou2023ifeval}
multi-constraint generator with nine independent, parameter-carrying verifiers
(e.g.\ \texttt{exact\_word\_count}, \texttt{forbidden\_word},
\texttt{word\_max\_length}, \texttt{num\_bullets}). Joint satisfiability of every
generated task is guaranteed by a witness the generator validates to score
$1.0$ (verified exhaustively: $3{,}189$ easy-tier and $19{,}632$ hard-tier
combinations). The \emph{hard} tier mandates \texttt{exact\_word\_count} (a
genuine LLM skill gap) plus at least one further ``brutal'' constraint, which
structurally caps the seed below saturation.

\paragraph{Metric.} Mean fraction of constraints satisfied on a disjoint
held-out split (\texttt{mc\_score} $\in [0,1]$), scored programmatically (no LLM
judge). We report a relative result under matched budget, not a
leaderboard number.

\paragraph{Reproducibility and budget.}
Every LLM call is recorded to a content-addressed cassette (\texttt{sha256} of the call's name,
positional arguments \emph{and} keyword arguments, which is what gives a per-replicate salt or a
changed temperature its own key rather than the previous draw's), so
runs replay offline at zero cost and partial spend is always persisted. Reproduction is strict:
re-running any study with \texttt{-{}-replay} re-executes the whole pipeline (loaders, GEPA search,
scoring) against the cached calls with no API key, and any cassette miss raises, per call and,
after the correction recorded as defect~7 of Paper~III's defect table, per run as well, since a
consumer that swallows the per-call
error can otherwise let a replay finish with a different number. GEPA is deterministic under
replay, so replay is bit-for-bit, and \texttt{scripts/verify\_reproducibility.py}
registers thirty-eight committed runs. Thirteen are labelled \textbf{superseded} and each
carries its own reason: eleven because the corrected harness changes every reflection prompt and
therefore every cassette key, and two because replaying them re-serves an empty reflection
completion that the blank-gene guard now refuses. The script reports them as skipped rather than
verified.
The remaining twenty-five were recorded under the corrected harness and every per-arm held-out score replays
exactly. The five remaining runs of the forty-three released are the
BigCodeBench ones \citep{zhuo2024bigcodebench}, listed as unreplayable because executing model code re-rolls unseeded
randomness and so shifts the content-addressed prompt key. Those three states partition the
release with nothing left over, and a test asserts they do: a run that supplies a number in
these pages cannot hold no reproducibility status at all. Output tokens are hard-capped via enforced \texttt{max\_tokens}; a pre-flight guard
refuses any run whose conservative estimate exceeds \texttt{-{}-max-usd}, whose default is \$5.
Five runs raised it explicitly, at four distinct values (\$6.18 twice, \$5.89, \$6.52, \$11.76), so
\$5 is a default rather than a study-wide ceiling. Experiments were pre-registered: hypothesis,
frozen configuration, expected seed band and verdict rule committed to version control
\emph{before} spending. Most of them, not all, and the exceptions are named here rather than
left for a reader to reconstruct from commit dates. The first single-constraint run predates the practice and
has no registration. The code-domain runs and the actuator experiment have none either, each
appearing first in the commit that carries its own results. Two further lines, the
feedback-actionability pair and the second-worker replication, do carry registrations, but
committed in the same commit as their results and after the runs' own recorded timestamps, by
seventeen to twenty-five minutes in the first case and by two and a half to three and a quarter
hours in the second. Those registrations may well have been written first; version control cannot
show it, so for them we do not claim it.


\paragraph{Four arms.}
All arms share an identical seed genome; they differ only in what GEPA may
evolve, so a held-out gap isolates componentized optimization from a
``more prompt text'' confound.
\begin{itemize}
  \item \textbf{A} --- seed baseline, evolves nothing.
  \item \textbf{B} --- prompt-only: evolves the \texttt{generate} gene.
  \item \textbf{C} --- command control: evolves the whole genome but scores with
    a genome-blind worker; its landscape is flat by construction (a harness
    soundness check).
  \item \textbf{D} --- loop-level: evolves the whole genome
    (\texttt{generate}\,$+$\,\texttt{ensure\_correct}, and \texttt{plan} when
    enabled) with the LLM worker.
\end{itemize}
$\dd = \mathrm{mean}(D) - \mathrm{mean}(B)$ is the quantity of interest.
A and C are expected flat. \textbf{Caveat, established after the fact} (Paper~III): as originally
implemented neither control could fail --- A copied its
own seed score with $\delta$ hardcoded to $0.0$, and C's command loop scored a saturated
$1.0$, so it proposed nothing. The corrected harness re-measures A and gives C a
genome-blind constant draft scored by the real metric, which is flat without being
saturated.

\paragraph{Optimizer.}
We drive the official GEPA \citep{agrawal2025gepa} via a loop-genome adapter, calling
\texttt{gepa.optimize} with \texttt{seed\_candidate}, \texttt{trainset}, \texttt{adapter},
\texttt{reflection\_lm} and \texttt{max\_metric\_calls}, plus a
\texttt{MaxCandidateProposalsStopper} to make the reflection-LM spend a hard cap. GEPA's reflection channel (offline, between
training iterations) sees per-constraint feedback prose --- this is the standard
optimizer signal and is identical across B and D; it is distinct from the
in-loop critique gate.

\paragraph{Two levers.}
(i) \emph{Feedback richness}: the critique step is optionally given
\texttt{text\_stats(draft)}, a deterministic one-line readout of measured
properties of the draft (word/character/sentence/bullet counts, longest-word
length, last words). It takes only the text and never sees the task, so it is a
\emph{measurement}, the tool/executor signal \citep{huang2023cannot} found
effective --- not the oracle. (ii) \emph{Complexity}: an optional third
\texttt{plan} gene turns the loop into plan--generate--critique--refine, giving
D a third interacting module to co-tune (the DSPy lever); B still evolves only
\texttt{generate}. Paper~II reports what these levers buy the loop itself; here they are
context for the contrast.

\paragraph{Verdict rule (pre-registered).}
\textbf{SUPPORTED} iff $\dd > \sigma_B + \sigma_D$ and positive;
\textbf{NOT SUPPORTED} iff $\mathrm{mean}(D) \le \mathrm{mean}(B)$;
\textbf{INCONCLUSIVE} if positive-but-within-noise, fewer than two seeds, or a
ceiling/floor flag fires (seed held-out $\ge 0.9$ or $\le 0.1$). This is a
deliberately conservative bar (the sum of two standard deviations), erring
toward under-claiming.

\section{Experimental Setup}
\label{sec:setup}

\paragraph{Models.} Worker \texttt{gpt-4.1-mini} (output cap 320 tokens); reflection
\texttt{gpt-5.6-luna} on every run recorded under the corrected harness --- the four-arm
re-run, the $p_0$ sweep and the split-size sweep --- and \texttt{gpt-5.6-terra}, which is the
CLI default, on the superseded runs that preceded it (output cap 1500 throughout). The
reflection model is therefore \emph{not} constant across the superseded and corrected sets, so
it is one more reason not to read a superseded number against a corrected one; within each set
it is fixed, which is what the contrasts require.

\section{Results}
\label{sec:results}

We report the four-arm contrast first (\S\ref{sec:rerun}), then the two subsections that change
how it should be read: \S\ref{sec:p0} varies $p_0$ and finds the contrast's sign is conditional
on it, and \S\ref{sec:decomp} separates the part of that contrast which is actually the
treatment from the part contributed by the gene both arms share. \S\ref{sec:analysis} gives the
mechanism, \S\ref{sec:posctrl} the positive control, and two pre-registered sweeps
(\S\ref{sec:prereg}, \S\ref{sec:budgetprereg}) close the alternative readings. A cross-domain
replication on frontier-hard code follows (\S\ref{sec:bcb}). Fifteen earlier contrast runs are
superseded and reported, with the reason, in Paper~III.

\subsection{The four-arm re-run on the corrected harness}
\label{sec:rerun}

The one four-arm contrast in this paper that measures what the study set out to measure was bought on a rebuilt harness, after the audit of Paper~III and under the prediction registered in \S\ref{sec:directions} and reproduced verbatim there. Same task family, budget and worker as run~3
(multi-hard, $\mathrm{mmc}=28$, 8 train / 10 held-out, one refine round, three seeds), with
the three corrections in place: a component-specific reflective dataset, an explicit
\texttt{module\_selector=all}, and a non-saturating arm~C.

We bought it twice: a first attempt was discarded because an empty reflection completion was
accepted as the verify gene, so the arm carrying the treatment ran with no critique at all
(Paper~III, defect~8). Both runs are released; everything below is the clean
re-run.

\textbf{The treatment was applied.} Arm~D's \texttt{genes\_changed} lists
\texttt{ensure\_correct} on \textbf{three of three} seeds, against $0/5$ for arm~D on the
superseded flagship run and $0/15$ pooling all its arms (Paper~III). The genes GEPA wrote are substantial detector rubrics of $3235$,
$3022$ and $2920$ characters against a $127$-character seed, and no completion anywhere in
the run was empty (shortest reflection output: $968$ characters). The selector fix shows up
alongside it: arm~D now receives proposal rounds at arm~B's rate rather than half of it,
which was the starvation of Evidence~3.

\begin{table}[h]
\centering
\begin{tabular}{llccc}
\toprule
Arm & GEPA evolves & seed & optimized & $\Delta$ \\
\midrule
A & nothing (baseline) & 0.788 & 0.788 & $+0.000$ \\
C & genome ignored (control) & 0.788 & 0.788 & $+0.000$ \\
B & \texttt{generate} & 0.788 & 0.788 & $\mathbf{+0.000}$ \\
D & \texttt{generate} $+$ \texttt{ensure\_correct} & 0.788 & 0.747 & $\mathbf{-0.042}$ \\
\bottomrule
\end{tabular}
\caption{The corrected-harness re-run, held-out means over three seeds. Prompt-only
optimization is a wash; arm~D comes out $0.042$ lower, which carries no directional claim --- it
is $1.44\sigma$ of measurement noise and three other training-split sizes at the same $p_0$ come
out positive (\S\ref{sec:splitresult}). Note what the controls do and do not establish: arm~A re-scores the seed genome,
which under a deterministic replay path must reproduce its first measurement, so it is a
determinism check rather than a flatness control; and arm~C's worker ignores the genome, so no
proposal can improve its training score and its flatness is true by construction. Critically,
\textbf{neither control bounds measurement noise}: both are deterministic replays of a fixed
genome, so their flatness says nothing about the variance of a \emph{fresh} draw. What does vary
is measured in \S\ref{sec:threats}: $\sigma(\dd) = 0.0291$.}
\label{tab:rerun}
\end{table}

Per-seed $\dd(D-B) = [-0.095, -0.080, +0.050]$, mean $\mathbf{-0.042}$, exact-resample 95\%
CI $[-0.095, +0.050]$, sign $p=1.00$. Prediction (ii) --- that $\dd$ would stay inside our
$\pm0.05$ SESOI (smallest effect size of interest, in the sense of
\citet{lakens2018equivalence}) --- holds on the point estimate, and only just; the interval does
not, so this is not an equivalence claim either, and at $n=3$ the sign test cannot fall below
$p=0.25$ regardless of the data. What distinguishes this null from the fifteen it replaces is
not its width but its content: the arms produced different genomes and different scores, so
the interval measures a contrast that exists.

The result we did not predict is Table~\ref{tab:rerun}'s last column. The two optimizing arms did
\emph{not} behave alike:
\textbf{prompt-only optimization is a wash} ($+0.025$, $0.000$, $-0.025$; mean exactly $0.000$) while \textbf{loop-level optimization costs} $-0.042$ ($-0.070$, $-0.080$, $+0.025$). We read that, at the time, as small-split overfitting: eight training tasks against ten
held-out, with arm~D carrying two genes of surface to overfit rather than one. Three later results
dismantle the reading, and the paragraph is left standing so the correction is legible rather than
quietly absorbed. \S\ref{sec:decomp} attributes $-0.025$ of the $-0.042$ to the gene both arms
evolve; the pre-registered split sweep (\S\ref{sec:splitresult}) finds no split-size trend to
overfit along, a \emph{smaller} split giving $+0.017$ and larger ones $+0.062$ and $+0.023$ at the
same $p_0$ and the same held-out tasks; and the noise floor (\S\ref{sec:threats}) puts this cell at
$1.44\sigma$. So the honest summary is neither that loop-level optimization fails to help nor that
it pays a price: this is the one configuration of four where it came out negative, by an amount
the study cannot separate from measurement noise. The degradation runs 4 and 5 showed --- on the
pre-correction harness, so read as prompt-only behaviour rather than as a contrast (Paper~III) --- and the
reported instability of prompt optimizers on some tasks
\citep{yang2023opro,zhou2025bbopromptscale}, remain the relevant context; what they no longer
support is a directional claim from this cell.

Arm~B is exactly flat on this clean re-run, so no general claim that optimizing at all costs
held-out score is available here; the loss is specific to the arm carrying the extra gene. The
contaminated first attempt supported the broader reading, which
Appendix~\ref{sec:corrections} withdraws.

\paragraph{Where the contrast lives.}
Because every run now releases per-task held-out vectors, the seed-level means can be opened
up. Arms~B and~D score the same ten held-out tasks per seed, so the contrast decomposes into
task flips: $5/10$ tasks differ on seed~0 (B better on 4, D on 1), $4/10$ on seed~1 (B better
on 3, D on 1), and $2/10$ on seed~2 (D better on 2, B on 0). Pooled, $19$ of $30$ task scores
are exact ties and the entire $\dd$ rests on eleven flips, $7$ favouring B and $4$ favouring
D. We report this as description, not as extra power: the tasks
are nested within seeds, since one training split produces the genome pair that scores all
ten, so the inferential unit remains the seed and $n$ is still $3$.

Arm~C is worth one qualification, because its number is unchanged from the retracted runs and
what it demonstrates is not. Previously it scored a saturated $1.0$, so GEPA saw a perfect
program and proposed \emph{nothing}: no part of the search ran. Here it starts at $0.23$ to
$0.36$, GEPA proposes candidates, evaluates them, and accepts none on any seed. We do not report
a per-arm proposal count, because the released cassette is content-addressed and does not
attribute a reflection call to the arm that issued it; an earlier draft printed one, and it was
not derivable from the artifact. What the artifact does show is run-level and makes the same
point: the corrected re-run issued $32$ reflection calls over three seeds against the superseded
flagship's $22$ over five, so the search ran harder on fewer seeds. The
flatness
is still true by construction --- a genome-blind worker's training score cannot improve, so
acceptance is impossible --- but the control now exercises the machinery it is supposed to
check instead of exiting before it starts.

\subsection{The \texorpdfstring{$p_0$}{p0} sweep: when loop-level optimization pays}
\label{sec:p0}

\begin{figure}[!tb]
\centering
\includegraphics[width=\textwidth]{figures/fig-p0.pdf}
\caption{\textbf{The paper's answer, and the mechanism under it.} \emph{Top:} the loop-level
contrast $\dd(D-B)$ against $p_0$, the base rate of first drafts that already satisfy every
requirement. Open squares are the single draw each level was first published from; filled circles
are the five-replicate means of \S\ref{sec:repeat} with $95\%$ intervals on $t(4)$, which is the
distribution four degrees of freedom actually have. Small open circles are the three per-seed
values; the grey band is the measured noise floor $\pm\sigma(\dd) = 0.0291$ and the dashed rules
the pre-registered $\pm0.05$ SESOI. No endpoint is individually significant at $p<0.05$, and only $L_2$ reaches $p<0.10$: what
carries the finding is the $0.102$ swing, about $2.5\sigma$ once the denominator is allowed to
carry $\sigma$ twice. \emph{Bottom:} why the sign moves. As $p_0$ rises the evolved critique begins
to accept and is never wrong when it does, while the seed rubric accepts more often and is wrong
six, six and two times --- precision-$1$ acceptance, not looser acceptance. Every number is derived
from a released cassette by \texttt{scripts/make\_figures.py}, and
\texttt{tests/test\_figures.py} fails if any disagrees with the text.}
\label{fig:p0}
\end{figure}

Every result above was measured at one point on one axis. \S\ref{sec:analysis} explains the
loss with a mechanism that makes a quantitative commitment, and that commitment turns out to be
the difference between a null and a law. Write $p_0$ for the base rate of first drafts that
already satisfy every constraint. The training-signal \emph{cost} of an always-\texttt{FIX}
critique is the chance it damages a draft that needed no repair, which is proportional to
$p_0$. So the degenerate candidate is cheap when $p_0$ is small and expensive when it is large,
and the mechanism predicts a \textbf{crossing} rather than a constant:

\begin{quote}
Reflective optimization of a verifier gene converges on a non-accepting detector when $p_0$ is
low enough that unconditional refinement costs less on the training signal than discriminating
correctly; as $p_0$ rises, the same search must learn to emit \texttt{PASS}, and the loop-level
arm should stop losing.
\end{quote}

\paragraph{The lever, and why it was already in the data.}
\texttt{exact\_word\_count} is satisfied on only $0.16$ of held-out final drafts (Paper~II) and
the generator puts it in \emph{every} task of the hard tier, so
full-task success cannot exceed $0.16$ there: the tier's $p_0$ is pinned low by one mechanical
skill. Excluding it from the generated pool raises $p_0$ while leaving a brutal anchor in every
task. $p_0$ itself needs no dedicated measurement run --- the first critique event judged the
first draft, so any recorded cassette contains it
(\texttt{scripts/measure\_p0.py}) --- and one \$0.09 arm-A pass creates the cassette per level.
Three levels, arms~B and~D, three seeds, everything else held at \S\ref{sec:rerun}'s
configuration, \$1.79 per optimizing level.

\begin{table}[h]
\centering
\footnotesize
\begin{tabular}{clcccccc}
\toprule
Level & excluded & $\mathbb{E}[p_0]$ & $p_0$ & B seed$\to$opt & D seed$\to$opt & $\dd(D-B)$ & 95\% CI \\
\midrule
$L_1$ & none & 0.127 & 0.067 & $0.788 \to 0.788$ & $0.788 \to 0.747$ & $\mathbf{-0.042}$ & $[-0.095,+0.050]$ \\
$L_2$ & \texttt{exact\_word\_count} & 0.725 & 0.267 & $0.855 \to 0.918$ & $0.855 \to 0.978$ & $\mathbf{+0.060}$ & $[+0.005,+0.105]$ \\
$L_3$ & $+$ \texttt{word\_max\_length} & 0.832 & 0.667 & $0.942 \to 0.972$ & $0.942 \to 0.970$ & $-0.002$ & $[-0.045,+0.040]$ \\
\bottomrule
\end{tabular}
\caption{The $p_0$ sweep. $L_1$ \emph{is} the four-arm re-run of \S\ref{sec:rerun}, so it cost
nothing to include. The treatment reached \texttt{ensure\_correct} on $3/3$ seeds at every
level. \textbf{$\dd(D-B)$ swings by $0.102$ across the axis}, about $2.5\sigma$
(\S\ref{sec:threats}): $-0.042$ where almost no first draft is correct and $+0.060$ where a
quarter of them are, with all three seeds positive at $L_2$. The swing is the finding; the
$L_1$ endpoint is not separable from noise on its own and we no longer report it as a loss
(\S\ref{sec:splitresult}). $\mathbb{E}[p_0]$ is the independence estimate and it runs high by
roughly a factor of two, because constraint failures correlate --- at $L_1$ the measured $p_0$
is $2/30 = 0.067$ against $0.127$ predicted --- so the estimate orders the levels and the
measured column is what they are. $L_3$ is \emph{saturated}; its $-0.002$ is one draw, and five replicates put it at $+0.019$ (\S\ref{sec:repeat}), still not significant but not the zero this column shows.}
\label{tab:p0}
\end{table}

\paragraph{The registered predictions, and the one that failed.} We registered three before
paying. (i) \texttt{PASS} rate $0.00$ at $L_1$: \textbf{held}, on three of three seeds.
(ii) \texttt{PASS} rate $>0$ on at least two of three seeds at $L_3$: \textbf{held}, on exactly
two. (iii) $\dd(D-B)$ non-decreasing in $p_0$ across $L_1 < L_2 < L_3$: \textbf{falsified} ---
$L_3$'s $-0.002$ sits below $L_2$'s $+0.060$.

We report (iii) as failed rather than reinterpreted, and the reason is a design error of ours
that the reader should be able to hold against the result. At $L_3$ the seed genome already
scores $0.942$, leaving $0.058$ of headroom, and both arms consume nearly all of it ($0.972$ and
$0.970$). Excluding \texttt{word\_max\_length} on top of \texttt{exact\_word\_count} recreated
precisely the condition this paper already dismisses for the easy tier --- saturation, where no
contrast between arms can appear because there is no room for one. $L_3$ therefore cannot
discriminate the arms on the outcome measure, although its \emph{mechanism} measure still moved
as predicted. The informative comparison is $L_1$ against $L_2$, and a corrected high-$p_0$
level needs a difficulty knob that raises $p_0$ without raising the baseline --- which the
constraint-exclusion lever, being a difficulty knob, cannot do.

\paragraph{The mechanism moved, and in a sharper way than predicted.} Scoring each arm's
critique as a detector on the drafts it actually judged. Arm~B evolves only \texttt{generate},
so its critique \emph{is} the seed rubric and it serves as the comparator the earlier table
lacked:

\begin{center}\small
\begin{tabular}{llcccccc}
\toprule
level & critique & base rate judged & \texttt{PASS} emitted & tp & \textbf{fp} & fn & tn \\
\midrule
$L_1$ & B, seed rubric   & 0.033 & 6  & 0  & \textbf{6} & 1  & 23 \\
$L_1$ & D, evolved       & 0.067 & \textbf{0}  & 0  & \textbf{0} & 2  & 28 \\
$L_2$ & B, seed rubric   & 0.567 & 19 & 13 & \textbf{6} & 4  & 7 \\
$L_2$ & D, evolved       & 0.567 & 4  & 4  & \textbf{0} & 13 & 13 \\
$L_3$ & B, seed rubric   & 0.833 & 15 & 13 & \textbf{2} & 12 & 3 \\
$L_3$ & D, evolved       & 0.867 & 6  & 6  & \textbf{0} & 20 & 4 \\
\bottomrule
\end{tabular}
\end{center}

Read \emph{base rate judged} carefully: it is the share of already-correct drafts among the
ones that arm's own generator produced, which is not the level's $p_0$ (measured on the seed
genome, Table~\ref{tab:p0}) because both arms evolve \texttt{generate} too. At $L_1$ the two
coincide; at $L_2$ and $L_3$ both arms draft far better than the seed, so the drafts reaching
either critique pass more often than $p_0$ would suggest. Anyone recomputing a base rate from
$\mathrm{tp}+\mathrm{fn}$ in this table will get that column and not $p_0$, which is why it is
printed.

Three things follow. First, prediction (ii): arm~D's evolved critique stops refusing to accept
as $p_0$ rises, on $0/3$, $1/3$ and $2/3$ seeds. Second, and unpredicted, \textbf{it never
emits a false \texttt{PASS} at any level} --- pooled $\mathrm{fp} = 0$ throughout, against the
seed rubric's $6$, $6$ and $2$. So the search does not trade accuracy for acceptance; it
acquires \emph{precision-1} acceptance, accepting only where it is right and continuing to
refine everything else. Third, the comparator forces a correction to how we first stated this.
The evolved critique is \emph{more} conservative than the seed rubric at every level, not less:
at $L_2$ it emits four \texttt{PASS}es where the seed rubric emits nineteen, and it still wins
by $+0.060$. So the $L_2$ gain is not ``the loop-level arm learned to accept''. It is that
arm~D's critique buys its acceptances at precision~1 while arm~B's are $6/19$ wrong, and a
false \texttt{PASS} ends the loop on a draft that still needed repair. What changes with $p_0$
is how much that discrimination is worth.

\paragraph{What this does to the headline.} Loop-level optimization is not inert: its effect
depends on $p_0$, and the condition is the one the mechanism named. An earlier draft put this
more strongly --- \emph{conditionally harmful} --- and that word is withdrawn in
\S\ref{sec:splitresult}, because the low-$p_0$ cell it rested on is one configuration of four
and $1.44\sigma$ of measurement noise. What survives the withdrawal is the \emph{contrast}
between the levels, which is stronger than either cell alone and can be stated in the same
units. Treating the two levels as independent draws,
\[
\frac{\dd_{L_2} - \dd_{L_1}}{\sigma\sqrt{2}}
 \;=\; \frac{0.060 - (-0.042)}{0.0291\sqrt{2}} \;=\; \frac{0.102}{0.0412} \;\approx\; 2.5 ,
\]
and this is a lower bound on the ratio, since $\sigma$ was estimated on the unexcluded pool
while $L_2$'s easier tasks sit nearer the ceiling and will vary less. So what the $p_0$ axis buys
is not a loss turning into a gain; it is a swing of $0.102$ in $\dd$ that noise does not account
for, from no measurable effect to the study's largest loop-level gain. Every
instruction-following run in this paper, superseded or not, sits at the low-$p_0$ end of the axis
--- the hard tier mandates \texttt{exact\_word\_count}, which caps $p_0$ at $0.16$
(Paper~II) --- and where we measure it directly it is $0.067$: the end of the axis
where refusing to accept is nearly free. At $p_0 = 0.267$ the same harness, optimizer, budget and
seeds give $\dd(D-B) = +0.060$ with every seed agreeing. One caveat belongs with that number: at
$n=3$ with three same-signed deltas the exact-resample interval \emph{must} exclude zero, so
``excludes zero'' here restates sign-unanimity and the sign test is at its floor ($p = 0.25$). It
is suggestive, not significant. The load-bearing evidence is the \emph{registered} swing, which
does not depend on the low-$p_0$ cell being negative. Saturation bias runs the safe way here too:
arm~D is nearer the ceiling ($0.960$, $0.975$, exactly $1.000$), so a ceiling compresses its
advantage rather than manufacturing it.

\subsection{How much of \texorpdfstring{$\dd$}{Delta} is the treatment?}
\label{sec:decomp}

Every $\dd(D-B)$ above is a \emph{whole-system} contrast, and the arms do not differ in one
component. Arm~B evolves \texttt{generate}; arm~D evolves \texttt{generate} \textbf{and}
\texttt{ensure\_correct}. Only the second is the treatment, so a gap between the arms mixes a
differently-evolved shared gene with the loop built on top of it, and $\dd$ alone cannot say
which moved. Writing $d$ for an arm's mean first-draft score and $f$ for its mean final score,
the split is an identity:
\[
\dd \;=\; f_D - f_B \;=\; \underbrace{(d_D - d_B)}_{\text{draft term: }\texttt{generate}}
\;+\; \underbrace{(f_D - d_D) - (f_B - d_B)}_{\text{loop term: the treatment}}
\]
Both terms come from the released artifact at \$0 --- the drafts by replaying each arm's evolved
genome against its cassette, the finals from the per-task vectors every run now emits
(\texttt{scripts/analyze\_delta\_decomposition.py}). The pooled $\dd$ it reconstructs agrees
with the arm means the paid runs reported, which is what makes it a decomposition of the
headline rather than of a different quantity.

\begin{table}[h]
\centering
\small
\begin{tabular}{lcccc}
\toprule
level & $\dd(D-B)$ & draft term (\texttt{generate}) & loop term (treatment) & treated share \\
\midrule
$L_1$ & $-0.042$ & $\mathbf{-0.025}$ & $-0.017$ & $40\%$ \\
$L_2$ & $+0.060$ & $\mathbf{+0.000}$ & $\mathbf{+0.060}$ & $\mathbf{100\%}$ \\
$L_3$ & $-0.002$ & $+0.005$ & $-0.007$ & --- \\
\bottomrule
\end{tabular}
\caption{$\dd$ split into the gene both arms evolve and the gene only arm~D evolves. The
treated share is omitted at $L_3$, where $\dd$ is too near zero for a share of it to mean
anything. Two readings change. Arm~D's $-0.042$ at $L_1$ \textbf{overstates the loop-level
penalty by a factor of $2.5$}: the majority of it is a \texttt{generate} gene that came out
worse, which is a fact about prompt search and not about optimizing a verifier. And the
$+0.060$ at $L_2$ is \textbf{entirely} the treatment, with the shared gene contributing
exactly nothing --- both arms' drafts score $0.870$ there.}
\label{tab:decomp}
\end{table}

The registered sign flip therefore survives the decomposition and is cleaner on the treated
component alone: the loop term runs $-0.017 \to +0.060 \to -0.007$. We report the aggregate
$\dd$ as the headline because it is the pre-registered quantity, but the loop term is the one
that answers the paper's title question, and a componentized-optimization result should print
both --- knowing that the optimizer \emph{changed} a component does not tell you that the
component is what moved the score.

\paragraph{Where the loss sits, and the mechanism it does not support.} The always-\texttt{FIX}
account says the cost of never accepting is borne by drafts that needed no repair, a stratum
whose weight is $p_0$. That is directly checkable. Partition held-out tasks by whether the
\emph{seed} genome's first draft was already correct --- a pre-treatment covariate, identical
for both arms, so conditioning on it cannot condition on a collider --- and read the strata
separately.

\begin{center}\small
\begin{tabular}{llccccc}
\toprule
level & stratum & $n$ & B & D & D$-$B & contribution to $\dd$ \\
\midrule
$L_1$ & already correct & 2  & 0.7750 & 0.8000 & $+0.0250$ & $\mathbf{+0.002}$ \\
$L_1$ & already wrong   & 28 & 0.7893 & 0.7429 & $-0.0464$ & $\mathbf{-0.043}$ \\
$L_2$ & already correct & 8  & 0.8875 & 1.0000 & $+0.1125$ & $+0.030$ \\
$L_2$ & already wrong   & 22 & 0.9295 & 0.9705 & $+0.0409$ & $+0.030$ \\
$L_3$ & already correct & 20 & 0.9775 & 0.9675 & $-0.0100$ & $-0.007$ \\
$L_3$ & already wrong   & 10 & 0.9600 & 0.9750 & $+0.0150$ & $+0.005$ \\
\bottomrule
\end{tabular}
\end{center}

At $L_1$ the already-correct stratum contributes $+0.002$: arm~D scores \emph{better} there
than prompt-only, on the very two drafts its critique rejected. The whole $-0.042$ sits on the
$28$ drafts that were already wrong. The stratum the mechanism names is also, at $p_0 = 0.067$,
structurally too thin to carry the result --- its weight is $2/30$, so with per-task scores in
$[0,1]$ its contribution is bounded by $0.067$ in magnitude whatever happens inside it. We
therefore \textbf{withdraw the outcome-side half of the always-\texttt{FIX} account}: the
training-signal argument for why a search converges on never-accepting stands and is what
predicted the sign, but ``it loses by damaging correct drafts'' is not what the $L_1$ loss
consists of. At $n=2$ this rules the channel out as the dominant term rather than estimating
it, and that limitation is the finding: a mechanism whose stratum has two tasks in it cannot be
confirmed at low $p_0$, only bounded.

What the loop term at $L_1$ does consist of is visible in the refine rounds. Arm~D refines
$30/30$ held-out tasks against arm~B's $24/30$ --- its critique never accepts, arm~B's falsely
accepts six --- yet each of arm~D's refine rounds is worth less, $+0.043$ against $+0.075$. The
refine step is net \emph{positive} at every level on the seed genome ($+0.072$, $+0.160$,
$+0.042$), so this is not a destructive repair step being invoked too often; it is that the
${\sim}3000$-character evolved rubric produces less productive repair prose than the
$127$-character seed rubric. That connects the $p_0$ result to this study's strongest
established lever, feedback \emph{actionability} (Paper~II, worth $+0.063$): the
evolved gene moved the threshold to zero and made the feedback worse, and at $L_1$ the second
effect dominates. One caveat is load-bearing and cuts against reading the per-refine gap as an
estimate: arm~B refines a \emph{selected} subset, having skipped the six drafts its critique
waved through, so part of that gap is selection rather than feedback quality.

\section{Why the Lever Is Inert at Low \texorpdfstring{$p_0$}{p0}}
\label{sec:analysis}

\paragraph{What the optimizer actually wrote, scored as a detector.}
The re-run makes answerable a question a hand-written gene could only half-address. Below we
show that a verify gene written \emph{by hand} against our own findings moved only the
detector's threshold and never gave it discrimination; the obvious objection was that a
search over prose might find what one author could not. Arm~D's evolved genes are real
rubrics of ${\sim}3000$ characters, so we scored each one as a detector on exactly the drafts
it judged, against the seed rubric on the drafts \emph{it} judged (offline, \$0;
\texttt{scripts/analyze\_evolved\_gene.py}). A same-drafts contrast does not exist, because
the generator and critique genes moved together, but the question here is whether the evolved
detector discriminates at all.

\begin{center}\small
\begin{tabular}{lccccc}
\toprule
critique & PASS rate & tp & fp & fn & tn \\
\midrule
seed rubric, seed 0 & 0.20 & 0 & 2 & 1 & 7 \\
\textbf{evolved}, seed 0 & \textbf{0.00} & 0 & 0 & 0 & 10 \\
seed rubric, seed 1 & 0.20 & 0 & 2 & 0 & 8 \\
\textbf{evolved}, seed 1 & \textbf{0.00} & 0 & 0 & 0 & 10 \\
seed rubric, seed 2 & 0.30 & \textbf{1} & 2 & 0 & 7 \\
\textbf{evolved}, seed 2 & \textbf{0.00} & \textbf{0} & 0 & 2 & 8 \\
\bottomrule
\end{tabular}
\end{center}

\textbf{The evolved critique emits \texttt{PASS} zero times, on every seed.} Pooled, the
$127$-character seed rubric produces one true positive and the ${\sim}3000$-character evolved
rubrics produce none. Given a large budget of prose and a reflector that could finally see
its own verdicts, the optimizer converged on the same degenerate corner the hand-written gene
slid toward: never accept. This is the strongest form of the C1 finding, because it now holds
for the \emph{search} and not merely for one author. A critique that never passes sends every
draft to a refine round, including drafts that were already correct: on seed~2 two compliant
drafts were rejected. The loop-level arm does not lose because the optimizer failed at its
job; it loses because it succeeded at optimizing a training signal that rewards refusing to
accept. Structurally this is reward gaming in the sense of \citet{skalse2022reward}: the
training score is a proxy on which ``never accept'' is near-optimal at low $p_0$, and improving
the proxy is what moves the true objective the wrong way. Their result that unhackable proxies
are essentially unavailable is worth stating here, because it says the defence is not a better
training signal but a screen on the optimizer's output --- which is what a \texttt{PASS} rate of
zero provides. The same convergence has been reported independently, with the same optimizer and
a different degenerate corner: on the tasks where GEPA did change their prompt,
\citet{pappu2026multiagent} report it converging on \emph{extreme-deference} instructions telling
non-expert agents to treat the expert's answer as correct by default. Refusing to accept and
deferring unconditionally are the same move --- collapse the component to a constant policy that
the training signal cannot punish --- which suggests the corner is a property of reflective search
under a weak signal rather than of our task family.

One step of that account does not survive measurement, and \S\ref{sec:decomp} is where we
check it. The natural next sentence --- that refining an already-correct draft breaks it, and
that this is what the $-0.042$ consists of --- is wrong on both counts here. Those two
compliant drafts are the entire already-correct stratum at $L_1$, and arm~D scores $0.8000$ on
them against prompt-only's $0.7750$: the refine round did not break them. The deficit sits on
the drafts that were already \emph{wrong}. The training-signal argument for \emph{why the
search goes here} stands; the outcome-side argument for \emph{what it costs} needs the
decomposition, which also finds that $-0.025$ of the $-0.042$ is not this gene at all.

Why that corner is so attractive is not specific to our setup. On multi-hard almost every first
draft genuinely fails, so a detector that emits \texttt{FIX} unconditionally is nearly
indistinguishable \emph{on the training signal} from one that reasons correctly, and it is far
easier to express in prose; the penalty appears only on the minority of drafts that were already
compliant, which a small training split barely contains. Any componentized optimizer pointed at a
verifier on a task with a low base rate of correct first drafts faces the same gradient, which
makes ``never accept'' a failure mode to test for rather than a quirk of this run.

\paragraph{The same corner is reachable by hand, which is why the re-run was worth buying.}
Before paying for it we wrote a verify gene by hand against this paper's own findings and ran it
against the seed gene on one shared cassette, so each task's first draft was byte-identical across
the two conditions and the pair isolates the gene exactly. The hand-written gene moved the
detector's \emph{threshold} and nothing else: \texttt{PASS} rate $0.09 \to 0.00$, no true positive
on either side, held-out $-0.033$ with a bootstrap interval of $[-0.095,+0.028]$ and sign test
$p=0.61$ --- inconclusive under our registered rule, and not equivalence either. If the reachable
improvements in prose space are threshold moves, then a search over prose searches a space whose
gains are threshold moves, and a threshold move in a detector that cannot discriminate has no
systematic gain to offer. That is a mechanism for the loop-level null which does not involve the
optimizer at all, and it is evidence rather than proof: one hand-written point does not establish
that no prose adds discrimination, which is exactly what made the re-run above worth its price
(registered in \S\ref{sec:directions}). Appendix~\ref{sec:handgene} reports the experiment, its
detector table and the decomposition of the $-0.033$.

\paragraph{Intrinsic self-correction has a low ceiling, and optimizing on top of it
costs.} Consistent with \citet{huang2023cannot}, run~4 shows the loop improving when given an
external tool but the reflective optimization on top adding noise rather than signal --- the
``self-correction can make it worse'' regime --- and its optimized held-out scores fall below
the seed, echoing OPRO's observation that optimized instructions can overfit a training set
\citep{yang2023opro}. One qualification belongs with that reading: run~4 predates the harness
correction, so its \emph{contrast} between arms is not a measurement (Paper~III). What survives is
the one-gene arm, whose search the selector confound does not touch, and it is that arm's
optimized score falling below its own seed. With small budgets and small splits, proposals that win on the training
subsample lose on held-out, and this hurts B and D alike. \S\ref{sec:rerun} measures the
effect on a harness where the two arms actually differ.

\section{A Positive Control: the Instrument Does Respond}
\label{sec:posctrl}

A null is worth reading only if the instrument that produced it can register a signal, and
after three silent failures we do not ask the reader to take that on trust. We built a
landscape in which evolving the verify gene is the \emph{only} path to a passing answer: the
generator always drafts \texttt{x\,x}, which fails the task's single constraint however its
prompt is worded; the critique catches the flaw only when the verify gene carries a sentinel;
and a caught flaw triggers a refine that passes. Arm~B therefore cannot escape the failing
draft no matter how well it optimizes, and arm~D can. Running the real four-arm study through
real GEPA on this landscape (no stubbed optimizer, no network, no spend) gives
A~$0.000$, B~$0.000$, C~$0.000$, D~$\mathbf{1.000}$, i.e.\ $\dd(D-B) = +1.000$. The harness
detects a loop-level advantage when one exists, so the near-zero $\dd$ we report on real tasks
is a property of those tasks and that optimizer, not of the measurement path
(\texttt{scripts/positive\_control.py}). Paper~III argues for this as a standing requirement
rather than a courtesy, since it is the one artifact here that would have failed loudly under
the original wiring.

\section{Pre-registered: Is the Loss a Small-Split Artifact?}
\label{sec:prereg}
\label{sec:splitsweep}

Two mechanisms now account for arm~D's $-0.042$ at $L_1$: the training signal that rewards
refusing to accept (\S\ref{sec:analysis}), and the \texttt{generate} gene that carries $-0.025$
of it (\S\ref{sec:decomp}). One alternative explanation survives both, and the $p_0$ result
does not touch it: eight training tasks is a small split, arm~D has two genes of surface to
overfit rather than one, and the reflective optimizer overfits. That is testable by varying the
split, the harness already exposes \texttt{-{}-n-train}, and we pre-register the design here,
before paying, in the same form as the prediction that Section~\ref{sec:rerun} resolved.

\paragraph{A fixed-budget split sweep is not a split sweep.} The obvious experiment --- raise
\texttt{-{}-n-train} and hold \texttt{-{}-max-metric-calls} fixed --- is confounded in the
direction that manufactures the result we are looking for. GEPA charges its seed evaluation
$|\mathrm{valset}|$ metric calls and our \texttt{optimize\_loop} passes the train split as the
trainset with no separate valset, so one accept cycle costs
$2n_{\mathrm{train}} + \min(3, n_{\mathrm{train}})$: at $n_{\mathrm{train}}{=}24$ that is $51$
calls, and run~3's budget of $28$ would have bought an arm that cannot complete a single cycle,
reported \texttt{genes\_changed=[]} and a flat $\dd$, and read as \emph{optimization did not help}
when it meant \emph{optimization never ran}. This is Evidence~3's confound in a new place; we
caught it before the money rather than after the write-up, and the harness now refuses the
configuration at the dry-run step with the arithmetic in the error text. The refusal is scoped to
arms that evolve something, so the two \texttt{-{}-arms A} verifier-ablation runs, where the budget
is inert, still stand.

The sweep therefore holds \emph{metric calls per training task} constant at $3.5$ --- a
constant number of epochs over the split --- which reproduces run~3's budget exactly at
$n_{\mathrm{train}}{=}8$ and scales it elsewhere. The cost of that choice is stated plainly:
split size and total budget move together, so the sweep identifies their joint effect along
the constant-epoch line and not either one alone. They cannot be separated below the
accept-cycle floor, because the cell that would separate them is the degenerate one the guard
rejects; a budget-only sweep at fixed $n_{\mathrm{train}}{=}8$ (D4) is the complement.

\begin{table}[h]
\centering
\begin{tabular}{ccccc}
\toprule
$n_{\mathrm{train}}$ & \texttt{mmc} & accept-cycle floor & type-overlap count (3 seeds) & priced bound \\
\midrule
4  & 14 & 11 & 2 / 1 / 3 & \$1.47 \\
8  & 28 & 19 & 3 / 2 / 3 & (run~3) \\
16 & 56 & 35 & 5 / 5 / 3 & \$2.43 \\
24 & 84 & 51 & 6 / 6 / 4 & \$3.08 \\
\bottomrule
\end{tabular}
\caption{The pre-registered sweep. Arms~B and~D only: arm~A re-scores a fixed genome and
arm~C is flat by construction, and replaying run~3's cassette under \texttt{-{}-arms BD}
reproduces its B and D fields bit-for-bit, so dropping them changes nothing but the bill.
\emph{Accept-cycle floor} is $2n_{\mathrm{train}} + \min(3, n_{\mathrm{train}})$, the budget
below which the guard refuses. \emph{Type-overlap count} is how many of the
$n_{\mathrm{train}}$ training tasks have a constraint-type set that also occurs in held-out,
per seed; as a \emph{share} of the split it does not increase from the smallest split to the
largest on any seed ($2/4 \to 6/24$, $1/4 \to 6/24$, $3/4 \to 4/24$), so a larger split is not
leaking more. No exact task instance is shared at any size. Each cell
sits under the study's own \$5 ceiling; the bounds are the guard's upper estimates, and run~3
came in at $17\%$ of its own bound.}
\label{tab:splitsweep}
\end{table}

\paragraph{What is held fixed, verified offline at \$0.} The pool is
$\max(64,\,2\max(n_{\mathrm{train}}, n_{\mathrm{heldout}}) + 2)$, hence $64$ for every
$n_{\mathrm{train}} \le 31$ and larger from $32$. Because the pool size seeds the shuffle,
crossing that cliff silently changes \emph{which} tasks are held out, so the sweep stays below
it and a test pins the cliff. Within it we verify, for all three seeds: the ten held-out tasks
are identical at every sweep size; the training splits are nested prefixes, so the sweep adds
tasks rather than resampling them; train and held-out share no task id and no exact instance;
and the type-overlap shares of Table~\ref{tab:splitsweep}. These are properties of the
generator, not of a purchase, so they are asserted in the test suite rather than reported once.

\paragraph{The prediction.} The mechanism of \S\ref{sec:analysis} predicts that a larger split
will \emph{not} rescue arm~D. What the search converges on is a critique that never returns
\texttt{PASS}, and unconditional refinement is near-free on the training signal whenever most
first drafts fail --- a property of the task's pass rate, not of the split size. So the
degenerate candidate stays optimal on train at every $n_{\mathrm{train}}$ we can afford, and
we predict (i) the evolved critique's \texttt{PASS} rate remains $0.00$ on at least five of the
six arm-D genes at $n_{\mathrm{train}} \in \{16, 24\}$, and (ii) $\dd(D-B)$ at
$n_{\mathrm{train}}{=}24$ does not cross zero. Prediction (i) is the load-bearing one: it is a
near-deterministic property of a released gene, checkable offline for \$0 by
\texttt{scripts/analyze\_evolved\_gene.py} --- the script that produced the detector counts
earlier in this section --- whereas (ii) is a three-seed mean whose interval will
again be wide enough that a sign flip alone would settle nothing. The honest alternative is
live and would be the more useful paper: if \texttt{PASS} becomes non-zero and $\dd$ closes on
zero as the split grows, then the finding is not ``loop-level optimization is inert'' but
``loop-level optimization needs a training split larger than any we bought,'' and the headline
should change accordingly.

\subsection{Result: the load-bearing prediction holds and the headline does not replicate}
\label{sec:splitresult}

The three unbought rows were then bought, under the priced bound of \$6.98 stated above, and
Table~\ref{tab:splitresults} reports them beside run~3. The harness records the pre-flight
\emph{estimate} for each run and not the invoice, so what is verifiable in the artifact is the
bound ($\$1.47 + \$2.43 + \$3.08$) rather than the spend; run~3's precedent puts the estimator
at roughly $6\times$ conservative. The two predictions come apart, so we take them in the order
of their pre-registered weight.

\textbf{Prediction (i) holds, at exactly its stated threshold.} The evolved critique's
\texttt{PASS} rate is $0.00$ on $3/3$ arm-D genes at $n_{\mathrm{train}}{=}16$ and $2/3$ at
$n_{\mathrm{train}}{=}24$: five of six, which is what we pre-registered as the load-bearing
claim. Pooling all four sizes it is $11/12$. The single exception is instructive rather than
awkward: that gene emits exactly one \texttt{PASS}, and the draft it accepted was in fact
correct ($\mathrm{tp}{=}1$, $\mathrm{fp}{=}0$). Across every arm-D gene in the sweep the
evolved critique never emits a false \texttt{PASS}, which is the same precision-1 acceptance
reported at the $p_0$ levels. So the Goodhart mechanism of \S\ref{sec:analysis} is not a
small-split artifact: at low $p_0$ the reflective search converges on a non-accepting verifier
at every training-split size we can afford.

\textbf{Prediction (ii) is falsified on the point estimate.} We predicted $\dd(D-B)$ at
$n_{\mathrm{train}}{=}24$ would not cross zero; it reads $+0.023$. Two of three seeds are still
negative and the interval $[-0.025,+0.115]$ contains zero, so nothing is established in the
other direction either --- but the prediction was registered as a sign claim about a point
estimate, and as a sign claim about a point estimate it failed. We report it that way rather
than retreating to the interval, which is the move pre-registration exists to prevent.

\textbf{What the sweep actually establishes is about run~3, not about arm~D.} Every row of
Table~\ref{tab:splitresults} holds $p_0$, the held-out tasks and the seeds fixed; only the
split and its constant-epoch budget move. Across them $\dd(D-B)$ reads
$+0.017,\,-0.042,\,+0.062,\,+0.023$: three of four positive, and the one negative cell is
run~3 itself. The sign of the effect this paper leads with is therefore not a property of
$p_0 = 0.067$; it is a property of one $(n_{\mathrm{train}}, \mathrm{mmc})$ cell. We withdraw
the unconditional form of that claim in \S\ref{sec:threats} and state what survives:
\emph{the verifier reliably degenerates at low $p_0$, and we have not shown that this costs
held-out performance.} The two are separable and we had run them together.

The refine accounting says why the two can come apart, and it is visible in released data at
\$0. A never-accepting critique is behaviourally just \emph{always refine}, and at low $p_0$
refining is worth doing: arm~D refines $30/30$ held-out drafts at both $n_{\mathrm{train}}{=}4$
and $16$ and gains $+0.055$ and $+0.098$ per refined draft. A degenerate verifier and a harmful
verifier are different objects, and where almost every draft is wrong the degenerate one is
close to the right policy.

\begin{table}[h]
\centering
\small
\begin{tabular}{ccccccc}
\toprule
$n_{\mathrm{train}}$ & \texttt{mmc} & seed-genome & $\dd(D-B)$ & per-seed & $|\dd|/\sigma$ & arm~B \\
 & & held-out & & & & accepted \\
\midrule
4  & 14 & 0.7717 & $+0.017$ & $-0.020,\,+0.045,\,+0.025$ & 0.57 & \textbf{0/3} \\
8  & 28 & 0.7883 & $-0.042$ & $-0.095,\,-0.080,\,+0.050$ & 1.43 & 2/3 \\
16 & 56 & 0.7583 & $+0.062$ & $+0.055,\,+0.080,\,+0.050$ & 2.12 & 1/3 \\
24 & 84 & 0.7883 & $+0.023$ & $-0.020,\,-0.025,\,+0.115$ & 0.80 & 2/3 \\
\bottomrule
\end{tabular}
\caption{The sweep, bought. \emph{Seed-genome held-out} is the unoptimized genome's mean over
the same thirty held-out tasks in every row --- by construction one fixed quantity, measured
four times. It moves over a range of $0.030$, which is the measurement-noise estimate of
\S\ref{sec:threats} appearing inside the experiment that needed it; $\sigma = 0.0291$ in the
$|\dd|/\sigma$ column is that estimate. No cell reaches $2.5\sigma$ and the between-cell
spread of $\dd$ ($0.037$) barely exceeds the noise on a single cell, so these four numbers are
consistent with one quantity plus noise. \emph{Arm~B accepted} is the number of seeds on which
the prompt-only comparator changed a gene at all: at $n_{\mathrm{train}}{=}4$ it changed
nothing on any seed, so that row is not a $D$-vs-$B$ contrast but $D$ vs.\ the seed genome
(numerically identical here, and flagged rather than dropped --- see Evidence~9, Paper~III).
Excluding it, two of three rows are positive.}
\label{tab:splitresults}
\end{table}

\paragraph{The confound we priced, quantified after the fact.} Holding metric calls per
training task constant does not hold search depth constant. Promotable accept cycles,
$(\mathrm{mmc} - n_{\mathrm{train}})/(\min(3,n_{\mathrm{train}}) + n_{\mathrm{train}})$, rise
across the sweep as $1.43,\,1.82,\,2.11,\,2.22$, while the reflection minibatch falls from
$75\%$ of the split to $12\%$. Both move monotonically with $n_{\mathrm{train}}$ and neither
was held fixed, so a larger split also means a slightly deeper search over a less
representative minibatch. The starved arm~B at $n_{\mathrm{train}}{=}4$ is that confound
becoming visible: $1.43$ promotable cycles is enough budget to pass the guard and not enough
for a one-gene arm to reliably land an acceptance. This is the correct reading of the guard's
scope --- it enforces that an arm \emph{can} fund an accept cycle, never that one is accepted.

\section{Pre-registered: Is It a Search-Depth Artifact?}
\label{sec:budgetprereg}

The split sweep moved split size and metric-call budget together along a constant-epoch line, by
necessity (\S\ref{sec:splitsweep}), so it cannot say which of the two produced the swing from
$-0.042$ to $+0.062$. The complement holds the split fixed at eight and varies only the budget.
We pre-register it here, before paying, in the same form as the two predictions this paper has
already resolved.

\paragraph{A budget sweep that ignores the proposal cap measures nothing.} The obvious design ---
raise \texttt{-{}-max-metric-calls} and change nothing else --- is broken, and we caught it at the
dry-run step. GEPA's search here is bounded by \emph{two} independent limits: the metric-call
budget, which affords $(\mathrm{mmc} - n_{\mathrm{train}})/(\min(3,
n_{\mathrm{train}}) + n_{\mathrm{train}})$ promotable accept cycles, and
\texttt{-{}-reflection-calls}, a hard cap on proposals per arm wired through to GEPA's stopper.
Effective depth is the \emph{minimum} of the two. Every contrast run under the corrected harness
used a cap of $6$ and every one was budget-bound at $1.43$ to $2.22$ cycles, so the cap has never
bitten and no published number is affected --- but it saturates at
$\mathrm{mmc} = 8 + 6 \times 11 = 74$, so budgets of $80$ and $160$ buy the \emph{same} six
proposals and must return the same result. Our first draft of this sweep was
$\mathrm{mmc} \in \{20,40,80,160\}$ at the inherited cap, which would have spent \$7.93 of priced
bound on two cells that are one experiment and then reported a ``$12\times$ budget range'' whose
effective range was $1.09$ to $6$. That is the same confound in a third place, which makes it a
rule: \textbf{before sweeping a budget, check every other limit that could bind first.}

\begin{table}[h]
\centering
\small
\begin{tabular}{ccccc}
\toprule
\texttt{mmc} & cap & promotable cycles & effective depth & priced bound \\
\midrule
20  & 6  & 1.09  & 1.09 & \$1.65 \\
28  & 6  & 1.82  & 1.82 & (the re-run, free) \\
40  & 6  & 2.91  & 2.91 & \$2.01 \\
56  & 6  & 4.36  & 4.36 & \$2.29 \\
80  & 6  & 6.55  & \textbf{6.00} & \$2.72 \\
160 & 12 & 13.82 & \textbf{12.00} & \$4.94 \\
\bottomrule
\end{tabular}
\caption{The pre-registered budget sweep, at $n_{\mathrm{train}} = 8$ with held-out tasks, seeds
and every other lever fixed at the four-arm re-run's values. Arms B and D only: replaying that
run's cassette under \texttt{-{}-arms BD} reproduces its B and D fields bit-for-bit, which we
verified rather than assumed, so the $\mathrm{mmc}{=}28$ row is the four-arm re-run of \S\ref{sec:rerun} at no cost. \emph{Effective depth} is
$\min$(cycles, cap): the first five rows share the inherited cap of $6$ and so form a
single-knob budget axis, while the last row must raise the cap to exceed it --- that row moves
two knobs and is reported as such, because at a cap of $6$ no budget can buy more than six
proposals. Every cell sits under the study's own \$5 default ceiling, so unlike four earlier runs
this sweep needs no override.}
\label{tab:budgetsweep}
\end{table}

\paragraph{The predictions.} The mechanism of \S\ref{sec:analysis} says the degenerate verifier
is optimal on a training signal where almost nothing is already correct, which is a property of
$p_0$ and not of how long the search runs. So we predict (i) the evolved critique's \texttt{PASS}
rate remains $0.00$ on at least four of the six arm-D genes at the two deepest cells
($\mathrm{mmc}{=}80$ and $\mathrm{mmc}{=}160$). And we predict (ii) \textbf{no monotone trend in
$\dd(D-B)$ across effective depth}, with no cell exceeding $2.5$ times the directly measured
$\sigma(\dd) = 0.042$ at this level (\S\ref{sec:repeat}), i.e.\ $|\dd| < 0.105$ throughout. The
grounds: within the split sweep, depth already ranged $1.43$ to $2.22$ with no ordering
($+0.017$, $-0.042$, $+0.062$, $+0.023$), and five three-seed cells cannot resolve a small trend.

\paragraph{The verdict rule, and what would change our minds.} If $\dd$ at $12$ proposals exceeds
$+0.05$ (the SESOI) \emph{and} $2\sigma$, prediction (ii) fails and the honest headline becomes
\emph{loop-level optimization needs more search than we bought} --- the live alternative this
paper has flagged since the retraction, and the one outcome that would turn a null into a
positive result. If no cell clears $2\sigma$ and the ordering is not monotone, the under-search
explanation is rejected over an $11\times$ depth range and the null is that much more robust.
Anything between is reported as inconclusive rather than resolved.

\paragraph{Correction, added after the runs.} The depth arithmetic in this section is
\emph{wrong}, and \S\ref{sec:budgetresult} reports how. It charges every proposal a full valset
re-evaluation, which only an accepted proposal pays, so the cap binds from about
$\mathrm{mmc}=40$ and not from $74$. The registered text is left as it stood --- editing a
pre-registration after seeing the data is the thing pre-registration exists to prevent --- and the
consequence, three cells sharing one search depth, is reported as a cost rather than presented as
a design.

\paragraph{Two risks we state in advance.} The $\mathrm{mmc}{=}20$ cell affords $1.09$ promotable
cycles, i.e.\ about one proposal per arm, so either arm may accept nothing and produce a
one-sided or empty contrast. The budget guard permits it --- it enforces that an arm \emph{can}
fund an accept cycle, never that one is accepted (Evidence~9) --- so that row will be labelled by
the acceptance census and excluded from the trend test rather than read as a null. Second, these
five runs re-measure the same unoptimized genome on the same thirty held-out tasks as the four-arm
re-run, so they will nearly double the replicate count behind the cross-run $\sigma$ estimate
($5 \to 9$). That estimate will therefore move, and every ratio derived from it will be
re-derived; we note it now so the change reads as the improvement it is rather than as a
convenient adjustment.

\subsection{Result: not a search-depth artifact either}
\label{sec:budgetresult}

\begin{figure}[!tb]
\centering
\includegraphics[width=\textwidth]{figures/fig-depth.pdf}
\caption{\textbf{The last route to a positive result, closed.} \emph{Top:} $\dd(D-B)$ against
\emph{measured} search depth --- GEPA iterations counted by replaying each cassette, not the
registered formula this section corrects. Labels are each cell's \texttt{max\_metric\_calls}; open
markers are budget-bound cells and filled ones have reached the proposal cap. The band is the
directly measured within-run $\sigma(\dd) = 0.047$ and the dashed rule the $\pm0.05$ SESOI. Over a
six-fold depth range $\dd$ does not order along depth, and the deepest cell is the flattest of the
six at $+0.0017$, or $0.04\sigma$, against a registered rule that required it to clear both
$+0.05$ and $2\sigma$. \emph{Bottom:} the dissociation, and the sweep's most useful result.
Acceptance rises with budget while $\dd$ does not, so the optimizer is searching and accepting and
what it accepts does not transfer. The shaded column marks three cells that share one measured
depth: they are independent draws of a single configuration, which is why they are not joined by a
line that would turn their offset into a slope.}
\label{fig:depth}
\end{figure}

All five cells were bought. Both registered predictions hold, and the section that registered
them contains an arithmetic error worth reporting first, because it is the same class of error
this series is about and we made it after congratulating ourselves for catching it.

\paragraph{The registered depth formula was wrong, so three cells landed on one depth.} We priced
search depth as promotable accept cycles, $(\mathrm{mmc} - n_{\mathrm{train}})/(\min(3,
n_{\mathrm{train}}) + n_{\mathrm{train}})$, which charges every proposal a reflection minibatch
\emph{plus} a full valset re-evaluation. Only an \emph{accepted} proposal pays the second part: a
proposal rejected on the minibatch costs the minibatch alone. The formula therefore prices the
expensive case and understates how many iterations a budget buys, and the proposal cap of $6$
binds from about $\mathrm{mmc}=40$ rather than from the $\mathrm{mmc}=74$ we registered. Three
cells --- $40$, $56$ and $80$ --- consequently ran the same six proposals, which cost roughly
\$5 of priced bound for one experiment bought three times. We stopped computing depth and started
measuring it: GEPA logs one iteration per proposal round, so replaying each cassette offline
recovers the count exactly (\texttt{scripts/analyze\_search\_depth.py}), and the depth column of
Table~\ref{tab:budgetresult} is that measurement rather than our formula.

The waste bought something we would not have designed. Those three cells are \emph{independent}
draws of one configuration --- same seeds, same tasks, same cap, but freshly sampled reflections,
so their evolved genomes differ --- and they agree: $-0.008$, $-0.008$, $-0.018$, mean $-0.012$.
An unintentional triplicate at the depth where the cap binds. We do not read their tightness as
precision: three draws of a quantity whose single-draw $\sigma$ is $0.047$ (\S\ref{sec:repeat})
have a sample sd that is itself uncertain by roughly a factor of two.

\begin{table}[h]
\centering
\small
\begin{tabular}{cccccc}
\toprule
\texttt{mmc} & cap & priced depth & \textbf{measured depth} & binding limit & $\dd(D-B)$ \\
\midrule
20  & 6  & 1.09  & \textbf{2}  & budget        & $+0.028$ \\
28  & 6  & 1.82  & \textbf{4}  & budget        & $-0.042$ \\
40  & 6  & 2.91  & \textbf{6}  & proposal cap  & $-0.008$ \\
56  & 6  & 4.36  & \textbf{6}  & proposal cap  & $-0.008$ \\
80  & 6  & 6.55  & \textbf{6}  & proposal cap  & $-0.018$ \\
160 & 12 & 13.82 & \textbf{12} & proposal cap  & $+0.002$ \\
\bottomrule
\end{tabular}
\caption{The budget sweep, bought, at $n_{\mathrm{train}}=8$ with held-out tasks, seeds and every
other lever fixed. \emph{Priced depth} is the pre-registered formula; \emph{measured depth} is the
highest GEPA iteration the run actually reached, recovered by replaying its cassette offline. The
formula is wrong in a consistent direction because it charges rejected proposals for a valset pass
they never take, so the cap binds three cells that were meant to differ. Measured, the axis still
spans $2$ to $12$ --- a factor of six --- and $\dd$ does not order along it. The
$\mathrm{mmc}=28$ row is the four-arm re-run, included at no cost.}
\label{tab:budgetresult}
\end{table}

\paragraph{Prediction (ii) holds: no trend, and the deep cell is the flattest of all.} Across a
measured six-fold range of search depth, $\dd(D-B)$ reads
$+0.028,\,-0.042,\,-0.008,\,-0.008,\,-0.018,\,+0.002$. It does not increase with depth, it does
not order along depth, and no cell reaches one standard deviation of the directly measured
$\sigma(\dd)=0.047$ --- the largest is the original $-0.042$ at $0.89\sigma$. The registered
verdict rule required $\dd$ at the deepest cell to exceed both $+0.05$ and $2\sigma$ for the
under-search explanation to survive; it reads $+0.0017$ (the $+0.002$ of
Table~\ref{tab:budgetresult}, unrounded, because at this size the third decimal is the claim),
which is $0.04\sigma$. \textbf{So the
under-search explanation is rejected}: at this split and this task family, loop-level
optimization does not need more search than we bought, and the twelve-proposal cell is the
flattest in the table. That was the one remaining route to a positive result, and it is closed.

\paragraph{What budget does buy is acceptance, not performance.} Acceptance is a coarser readout
than $\dd$ and it moves where $\dd$ does not: counting arm-seed pairs where the optimizer changed
any gene, the sweep gives $4/6$, $5/6$, $4/6$, $5/6$, $6/6$, $6/6$ in depth order. The two
deepest cells are the only ones where every arm accepts on every seed. So more budget does make
the search accept more proposals --- it is not idling --- and those additional accepted proposals
do not move held-out score. That dissociation is the sweep's most useful result, because it
separates \emph{the optimizer failed to search} from \emph{the optimizer searched and found
nothing that transfers}, and it points at the second.

\paragraph{Prediction (i) holds, and exposes a limit on a claim we made.} We predicted the
evolved critique's \texttt{PASS} rate would stay at $0.00$ on at least four of the six arm-D
genes at the two deepest cells. It is $2/3$ at $\mathrm{mmc}=80$ and $2/3$ at
$\mathrm{mmc}=160$: four of six, the registered threshold exactly. The degenerate verifier is
therefore not an artifact of shallow search any more than of a small split. But the two
exceptions are not alike, and one of them qualifies \S\ref{sec:analysis}. At depth $12$ the
accepting gene emits a single \texttt{PASS} and is right ($\mathrm{tp}=1$, $\mathrm{fp}=0$),
consistent with the precision-1 acceptance reported at the $p_0$ levels. At depth $6$ the
accepting gene emits a single \texttt{PASS} and is \textbf{wrong} ($\mathrm{tp}=0$,
$\mathrm{fp}=1$) --- the first false \texttt{PASS} from an evolved gene anywhere in this study.
Our claim that the evolved critique ``never emits a false \texttt{PASS}'' was scoped to the $p_0$
levels and remains true there, but it does not generalize: precision-1 acceptance is a property
of those runs, not of the search.

\section{Cross-domain corroboration: BigCodeBench-Hard}
\label{sec:bcb}

To test whether the instruction-following null generalizes to a domain where the
verify step catches real \emph{semantic} errors and runs as an executable tool ---
the regime \citet{shinn2023reflexion} find self-correction works --- we replicated
the four-arm design on \textbf{BigCodeBench-Hard} \citep{zhuo2024bigcodebench}, hard
enough that strong models score far below their full-benchmark pass@1 on it (the Hard
subset is calibrated to stay challenging). An earlier attempt on MBPP+
\citep{liu2023evalplus} saturated its training split, and that failure is what motivated
the move to a frontier-hard benchmark.\footnote{Reported in Appendix~\ref{sec:actuator}.} Each task's \texttt{unittest} methods are partitioned into a visible half
(the in-loop, executable signal) and a disjoint hidden half (the held-out scorer),
reconstructing the anti-oracle split; solutions are graded in a pinned, offline Docker
sandbox (\texttt{-{}-network none}, memory/PID caps, single-thread BLAS). We ran arms
A--D on $2$ seeds, $8$ train / $10$ held-out tasks drawn from a sandbox-verified
$80$-task pool, worker \texttt{gpt-4.1-mini}, reflection \texttt{gpt-5.6-terra}, at
$\mathrm{mmc}=30$.

The seed baseline sits at $0.40$ (seed 0) and $0.50$ (seed 1) --- far from both floor
and ceiling, so the loop has genuine headroom, unlike MBPP+. And unlike MBPP+, GEPA
\emph{ran}: on seed 1 it accepted improving candidates for \emph{both} the prompt-only
arm B (evolving \texttt{generate}) and the loop-level arm D (evolving \texttt{generate}
\emph{and} \texttt{ensure\_correct}), each lifting the $8$-task training aggregate
$0.375 \to 0.625$; on seed 0 neither arm accepted a candidate within budget, and the
genome-blind control C is flat by construction. Yet on the held-out split \emph{every
arm scored exactly its seed baseline} ($0.40$, $0.50$): the accepted training
improvements changed which held-out tasks \texttt{gpt-4.1-mini} solves not at all. The
headline is therefore a clean $\dd(D-B)=0.00$ (exact-resample $95\%$ CI
$[0.000, 0.000]$): even when loop-level optimization has its extra lever and \emph{pulls} it,
an accepted \texttt{ensure\_correct} rewrite that improves training, train-split improvement
does not transfer and the lever is inert on held-out. The instruction-following null recurs in a
semantically-verified, executable-feedback domain.

\paragraph{Can GEPA \emph{discover} the actionable feedback?} No. Paper~II measures a $+0.20$ shift on this benchmark from
concrete, measurement-shaped repair feedback, hand-designed --- reported there without an interval,
because a seed on this benchmark re-partitions a fixed pool rather than drawing a fresh sample, so
read it as a large and consistent shift and not a significant one. Either way it is the one lever
known to move this domain, so the sharpest test of our thesis is whether the optimizer finds it when
pointed squarely at it. We added a fourth evolvable gene, an \texttt{actuator} that turns the visible-test
signal into the repair guidance handed to refine, and let arm~D co-evolve it with
\texttt{generate} while holding the detector fixed, so any gap isolates optimization of the
\emph{fixer's} feedback --- the one component we know can help. Across all three configurations
bought, $D \le B$, never $D > B$: a flat $0.00$ at $2$ seeds, and $\dd(D-B) = -0.10$ on a powered
replication where prompt-only did transfer. Appendix~\ref{sec:actuator} reports both runs and the
two corrections that apply to the paragraph as first written.

\paragraph{Reproducibility scope.} The main four-arm result reproduced identically across two
independent paid runs, and we release cassettes and scores. We do \emph{not} claim
byte-for-byte replay of the full call trace here: model solutions may use unseeded randomness,
so re-executing them can yield varying test output --- an intrinsic property of
executable-code benchmarks, not of our harness. (The instruction-following runs, which never
execute model code, do replay bit-for-bit.) We normalize the string-level non-determinism we
can: object ids, addresses, and temp paths in failure messages. The actuator runs share this
scope.

\section{Threats to Validity}
\label{sec:threats}

We state these plainly. \textbf{Power:} runs 4--10 use $n{=}3$ seeds (runs 2--3 use 5), and
our earlier claim that the intervals were ``unusually informative despite small $n$'' is
withdrawn as circular (Paper~III); we now report per-run tie counts alongside every
interval and an exact sign test. At $n{=}3$ that test cannot fall below $p=0.25$ even with
unanimous seeds, so no seed-level interval here is significant on replication grounds.
\textbf{Measurement noise, and
what it costs the headline:} no interval in this study contains within-condition sampling
variance, because the worker runs at the provider's default temperature with one sample per
task, so nothing was measured twice on purpose. It was nonetheless measured by accident, and we
report it rather than leaving the omission as a scope note. Ten released runs --- the four-arm
re-run, each row of the split sweep and each bought row of the budget sweep --- evaluate the
\emph{same} unoptimized seed genome on the \emph{same} thirty held-out tasks before optimizing
anything, so they are ten independent draws of one fixed number. Their means span $0.7400$ to
$0.8183$: a range of $0.078$ on a quantity that is constant by construction, which is the
plainest statement of the problem this section exists to quantify. We report the count and the
range rather than enumerating them, because any future run at this split size enrols itself as a
replicate and an enumeration would need editing every time.
From the paired differences, $\sigma$ for one task under one draw is $0.1125$; a $\dd$ compares
two independent draws, so $\sigma(\dd)$ on one task is $0.1591$ and on a thirty-task mean
\[
\sigma(\dd) \;=\; 0.1591/\sqrt{30} \;=\; 0.0291 .
\]
In those units the study's effects are: the $L_1$ loss $-0.042$ at $1.44\sigma$, the $L_2$ gain
$+0.060$ at $2.07\sigma$ and the $L_3$ null $-0.002$ at $0.06\sigma$. Paper~III shows this
single study-wide $\sigma$ is an approximation --- measured per genome it spans $5.9\times$ ---
and \S\ref{sec:repeat} supersedes it here with five-replicate standard errors wherever the two
overlap. \textbf{We therefore withdraw the unconditional claim that evolving
the verifier costs held-out performance at low $p_0$.} It rests on a single $1.44\sigma$ cell
that three other training-split sizes do not reproduce (\S\ref{sec:splitresult}). What survives
is narrower and better supported: the verifier reliably degenerates to non-acceptance at low
$p_0$ ($11/12$ genes across four split sizes), and the $+0.060$ gain at higher $p_0$ is the
study's largest loop-level effect and is unanimous across seeds. Two caveats on the estimate
itself: the mean signed difference
between replicates, $-0.014$, is confounded with any provider-side drift between recording
dates, and five replicates of three seeds cannot separate the two. It establishes an order of
magnitude, and that order of magnitude is the same as the effects being reported. Direct
replicate sampling (D7) is the fix, needs no re-optimization, and has since been run:
\S\ref{sec:repeat} reports it and supersedes this estimate wherever the two overlap.

\textbf{A do-nothing floor in the
metric:} the fraction-of-constraints-satisfied score is not anchored at zero for an empty
answer --- the empty string scores $0.548$ on the vendored IFBench slice and $0.420$ on
generated multi-hard, and five of eleven IFBench verifiers are vacuously satisfied by text
containing no whitespace, quotes or bullets (\texttt{\_sub\_bullets("")} reduces to
\texttt{all([]) = True}). These are faithful ports of the upstream checkers and the floor
applies to every arm equally, so it does not bias \dd, but it compresses the scale: on IFBench
the floor sits at or above our reported seed baselines. It also makes the metric a poor
instrument for this loop, since a metric that partly rewards saying less is exposed to a
critique that pushes toward shorter answers. That was open when first written, because no
released run carried its evolved gene text; the corrected harness releases all of them, so it
is now measured and did not happen. Mean final-answer length under the arm whose critique never
accepts is $27.3$ words against prompt-only's $26.4$ at $L_1$ ($42.3$ vs.\ $43.2$ at $L_2$,
$45.5$ vs.\ $46.1$ at $L_3$): the loop-level arm's answers are not systematically shorter, and
a refine round moves length by under two words at every level, against a floor that needs an
empty answer to reach. \textbf{Single worker/reflection
model:} results may be model-dependent \citep{huang2023cannot}. \textbf{Cross-lineage, not
cross-provider:} run 6's critic shares a vendor and training lineage with the worker, so it
breaks self-\emph{grading} but may retain self-preference bias (D3). \textbf{One task family:}
verifiable instruction-following may be a domain where the critique has little leverage.
\textbf{Budget:} $\mathrm{mmc}\in[20,64]$ was small; \S\ref{sec:budgetresult} now spans $20$
to $160$ and finds no trend in $\dd$. \textbf{Held-out generalization:} the
generator's parameter pools are small, so train and held-out tasks can share parameter
combinations; this affects all arms equally but softens the generalization claim.
Section~\ref{sec:directions} turns each threat into a concrete follow-up.

\subsection{Measured: 135 re-scorings of the three levels}
\label{sec:repeat}

The estimate above is an inference from runs that happened to re-measure one genome. We then
bought the direct version (D7): every released genome of the $p_0$ sweep re-scored $K=5$ times
on its own held-out tasks --- three levels $\times$ \{prompt-only, loop-level, unoptimized\}
$\times$ three seeds $\times$ five replicates, $135$ runs. No re-optimization is involved, so
this is the existing scoring path on a fixed genome, and all $135$ replay from their cassettes
offline. Two properties make it a measurement rather than a cache: repeated scorings of one
genome are \emph{identical} calls, so a per-replicate salt enters the content-addressed key
without reaching the provider (without it, record-mode dedup would hand $K$ replicates one
cached draw and report zero variance); and no temperature is pinned, because the point is the
noise in the numbers the paper already reports, which were drawn at the provider's default.

\begin{table}[h]
\centering
\small
\begin{tabular}{lcccccl}
\toprule
level & single draw & $\dd$ ($K{=}5$) & sd(draw) & se(mean) & $t$ & verdict \\
\midrule
$L_1$ & $-0.042$ & $-0.021$ & 0.0468 & 0.0209 & 1.00 & not significant \\
$L_2$ & $+0.060$ & $+0.055$ & 0.0459 & 0.0205 & \textbf{2.70} & $\mathbf{p < 0.10}$ \\
$L_3$ & $-0.002$ & $+0.019$ & 0.0220 & 0.0099 & 1.89 & not significant \\
\bottomrule
\end{tabular}
\caption{$\dd(D-B)$ at each $p_0$ level from five independent draws of the same released
genomes on the same held-out tasks. \emph{Single draw} is what the paper reports elsewhere.
$t$ is the five-replicate mean over its own standard error, on $4$ degrees of freedom, so the
critical values are $2.132$ ($p<0.10$) and $2.776$ ($p<0.05$) --- \emph{not} $1.96$. That
distinction decides the $L_2$ row, the one the paper leans on: a normal approximation would call
$t = 2.70$ significant at $p<0.05$ and $t(4)$ does not.
These standard errors describe sampling noise on a \emph{fixed} task set, since every replicate
scores the same thirty tasks; they contain no task-draw component and are not a substitute for
more seeds.}
\label{tab:repeat}
\end{table}

Three things follow, and they cut in different directions. First, \textbf{the withdrawal of the
$L_1$ loss is confirmed by direct measurement}: averaged over five draws it halves, from
$-0.042$ to $-0.021$, and sits $1.12$ standard errors from zero. Second, \textbf{the $L_2$ gain
survives, at $p<0.10$} --- $+0.055$ against its own five-replicate standard error, $t = 2.70$
against a $2.776$ threshold it misses by $0.08$. This is the only inferential statement in the
paper that does not rest on three seeds, and it is the study's best-supported loop-level effect.
It is not significant at $p<0.05$, and an earlier draft reported it as such on a standard error
built from a population sd (Appendix~\ref{sec:corrections}). Third, \textbf{$L_3$ is not the zero we
reported}. We described it as a saturated null on the strength of a single draw of $-0.002$;
five draws give $+0.019$ with $t=1.89$, which fails the $p<0.10$ threshold by $0.24$ and is
therefore still not a result --- but the sign was an artifact of one sample, and we say so
rather than leave the earlier reading standing.

\paragraph{Which threats the released artifact can close.} The threats differ in what each
requires: more spend against the same harness, a credential we do not hold, or new harness
code. Table~\ref{tab:closability} separates them, so a reader can tell which open questions
are a budget line and which are engineering.

One replication predates the audit. A worker swap to \texttt{gpt-4.1-nano} on the multi-hard
regime ($3$ seeds, $16$/$16$, $\mathrm{mmc}=64$) gives $\dd(D-B) = -0.014$ (corrected interval
$[-0.041,+0.000]$; two of three seeds are ties), arm~B at $0.70 \pm 0.03$ and arm~D at
$0.69 \pm 0.03$ from a seed of $0.67$. It shows the pattern is not an artifact of one worker
model, but it inherits the same harness, so it is not a test of loop-level optimization. An
$\mathrm{mmc}{=}160$ replication was designed, priced, started and stopped mid-run once the
audit landed, having spent $\approx\!\$0.33$ of a projected $\$0.90$; its partial cassette is
not released, since a recording with no corresponding result would be a misleading artifact.

\begin{table}[h]
\centering
\small
\begin{tabular}{p{0.27\textwidth}p{0.15\textwidth}p{0.44\textwidth}}
\toprule
Threat & Closable as released? & What it requires \\
\midrule
\textbf{Treatment not applied} (Paper~III) & \textbf{Closed} (\S\ref{sec:rerun}) &
  The four-arm re-run on the corrected harness changed \texttt{ensure\_correct} on $3/3$
  seeds. Every row below is now a threat to a measurement that happened \\
Single worker/reflection model & Partly, by spend & Worker swap done pre-audit
  ($\dd=-0.014$); needs redoing post-fix \\
\textbf{Small train split} ($8$) & \textbf{Closed} (\S\ref{sec:splitresult}) & Bought:
  $n_{\mathrm{train}} \in \{4,8,16,24\}$ at constant epochs, held-out and seeds fixed. The
  degenerate verifier survives every size ($11/12$ genes); $\dd(D-B)$ does not, and the
  $-0.042$ is the only negative cell --- which is why the loss claim is withdrawn rather than
  rescued \\
Budget, $\mathrm{mmc}\in[20,64]$ & \textbf{Closed} (\S\ref{sec:budgetresult}) & Bought to
  $\mathrm{mmc}{=}160$ under the corrected harness. Measured search depth spans $2$ to $12$
  iterations; $\dd$ does not order along it and the deepest cell is the flattest \\
Power, $n{=}3$; no sampling variance & Partly, as released & $\sigma(\dd) = 0.029$ is
  recoverable at \$0 from ten runs that re-score one fixed seed genome (\S\ref{sec:threats}),
  and it is already decisive: it withdrew a headline. Tightening it needs the repeat loop,
  which re-scores released genomes and needs no re-optimization (D7) \\
\textbf{Seeds are partitions, not samples} (Evidence~4b) & Partly, as released & Task-level
  pairing on the per-task scores runs now emit; for the recorded runs, a re-run \\
\textbf{Metric has a do-nothing floor} (\S\ref{sec:threats}) & Partly, as released &
  The exploitation it invites is now measured and absent (answer length by arm,
  \S\ref{sec:threats}); removing the floor itself needs verifiers that fail rather than
  vacuously pass on empty text \\
\textbf{$\dd$ credits the treatment with a shared gene} (\S\ref{sec:decomp}) & \textbf{Closed}
  (\S\ref{sec:decomp}) & Both arms evolve \texttt{generate}, so $\dd$ mixes it with the
  treatment; the draft/loop split separates them offline from the released artifact \\
Cross-lineage, not cross-provider & No --- credential & A Gemini or Anthropic key; pricing
  already in place (D3) \\
One task family (math) & No --- new harness & Parametric math generator plus numeric checker;
  instance-level visible/hidden split \\
Held-out generalization & No --- generator work & Enlarged parameter pools, then re-run and
  report the gap (D8) \\
\bottomrule
\end{tabular}
\caption{Threat closability against the released artifact and an OpenAI-only key. ``By
spend'' means the harness answers it with no new code. The first row was the one that
superseded all others --- a threat to the existence of a result rather than to its
interpretation --- and \S\ref{sec:rerun} closes it. The two remaining bolded rows came from
the second audit; neither questions whether the measurement happened, but both change what
its numbers may be read to mean.}
\label{tab:closability}
\end{table}

\section{When Would Loop-Level Help? Directions}
\label{sec:directions}

The null is not ``componentized optimization is useless''; it is ``the critique component of a
self-refine loop is not a productive optimization target when it is an unreliable,
self-referential detector on shallow-verification tasks.'' Five of the directions below were
bought during the study, and what they changed is the priority: no longer \emph{get a better
detector}, since a perfect one does not suffice, but \emph{fix the fixer} --- which Paper~II
shows to be partly feedback-limited and partly generator-limited. Each changes one variable.

\begin{description}\itemsep1pt
  \item[D1$'$ (fix-side interventions) --- part done.] Run 8 bought the refine-budget half
    ($R{:}1\!\to\!3$, $+0.065$). \emph{Remaining:} a stronger worker on refinement only, and
    constraint-specific repair \emph{tools} the model can act on rather than merely read.
  \item[D2$'$ (feedback actionability) --- \textbf{done}, runs 9--10 (Paper~II).]
    \emph{Remaining:} whether an LLM \emph{rewriting} the measured edits adds anything over the
    raw measurement.
  \item[D3 (true cross-family critic) --- \textbf{blocked on a credential}.] Run 6's critic
    shares a vendor and lineage with the worker. Repeating it cross-\emph{provider}
    (\texttt{gemini-\allowbreak 2.5-\allowbreak flash-\allowbreak lite} or
    \texttt{claude-\allowbreak haiku-\allowbreak 4-5}, pre-priced) at $\ge 8$ seeds is the only
    direction that separates an intrinsically unreliable detector from a self-preferring one.
  \item[D4 (search depth) --- \textbf{done, and closed} (\S\ref{sec:budgetresult}).] Bought at
    fixed $n_{\mathrm{train}}{=}8$ over a measured six-fold depth range: no trend, deepest cell
    flattest. That was the last configuration in which a positive finding was available, since
    the split sweep moved split size and budget together and could not separate them.
    \emph{Remaining:} worker $\in$\{nano, mini, a 5-series\}, which tests under-budgeting of the
    \emph{model} rather than of the search.
  \item[D5 (cross-domain) --- \textbf{done} (\S\ref{sec:bcb}).] The null recurs under
    executable feedback, so reliable, actionable verification is not by itself sufficient.
    \emph{Remaining:} math, i.e.\ a numeric checker.
  \item[D7 (statistics) --- \textbf{done}, $135$ re-scorings (\S\ref{sec:repeat}).] They settle
    all three $p_0$ effects, and Paper~III reports what they show about $\sigma$ itself.
    \emph{Remaining:} the same measurement in the code domain, whose $\sigma$
    column is unavailable rather than favourable. Two
    design notes for anyone repeating it: the salt must enter the cassette key, or record-mode
    dedup returns $K$ byte-identical draws and reports \emph{zero} variance; and a seed buys a
    fresh sample only where it does not merely re-partition a fixed pool (Evidence~4b), so the
    unit to pair on is the \emph{task}, which is why every run now releases its per-task vector.
  \item[D1$'$-tools (deterministic actuator).] Paper~II localizes the residual to
    \texttt{exact\_word\_count}, so a counter would recover most of it; its remaining value is
    confirming the other constraints survive enforcement.
  \item[D8 (generalization).] Enlarge the generator's parameter pools so train and held-out
    cannot near-duplicate, and report the gap. On the code side, a benchmark postdating the
    worker's cutoff \citep{jain2024livecodebench} removes the contamination confound
    \S\ref{sec:bcb} cannot, at the cost of a moving target a replayable study must pin.
  \item[D9 (a high-$p_0$ level that does not saturate) --- \emph{the open design problem}.] The
    falsified monotonicity prediction needs a knob that raises $p_0$ \emph{without} raising the
    baseline, and constraint exclusion cannot supply one: being a difficulty knob it moves both,
    and $L_3$ saturates. Two candidates separate them. (i) Hold the constraint mix fixed and
    vary the \emph{drafting} budget: a second draft-then-pick step raises the share of
    already-correct drafts entering the critique without making any task easier to score.
    (ii) Keep the hard tier and \emph{sample} tasks to a target $p_0$, now free to measure per
    task from any cassette, so the level is built to specification rather than discovered. Under
    either, the already-correct stratum stops being two tasks wide, which is what
    \S\ref{sec:decomp} needs to estimate the damage channel instead of bounding it.
\end{description}

\paragraph{A pre-registered prediction for the corrected-harness re-run --- \textbf{resolved,
\S\ref{sec:rerun}}.} Prediction (i) held decisively --- \texttt{ensure\_correct} changed on all
three seeds of arm~D --- and (ii) held on the point estimate only barely, $\dd(D-B) = -0.042$
against a $\pm0.05$ SESOI, with the interval $[-0.095,+0.050]$ extending well past it, so this is
not an equivalence claim. \emph{The pre-registration, as written:} the hand-written verify gene
of \S\ref{sec:analysis} changed the detector's threshold (PASS rate $0.09 \to
0.00$) without producing a true positive, and moved held-out by $-0.033$ ($p=0.61$). If that
generalizes to the search --- if the reachable improvements in prose space are threshold moves
--- then the corrected harness should produce a specific, unusual pattern: \textbf{the
treatment is applied and still does nothing}. Concretely we predict (i) \texttt{genes\_changed}
will now include \texttt{ensure\_correct} on at least one seed of arm~D, since the reflector can
finally see the detector's own verdicts and the selector no longer starves it; and (ii)
$\dd(D-B)$ will nonetheless remain within our $\pm0.05$ SESOI. Prediction (i) is the one that
matters procedurally: it is the check the released runs could not perform, and if it fails the
re-run is again not a test of loop-level optimization and must be reported as such rather than
as a null. Prediction (ii) is the substantive one, and the honest alternative is live --- an
optimizer explores a far larger region of prose space than one author, so it may find the
discrimination we could not write.

\section{Conclusion}
\label{sec:conclusion}

\textbf{Does loop-level optimization beat prompt-only?} \textbf{It depends on $p_0$, and that is
the paper's answer.} At $p_0 = 0.267$ it
gains $0.060$ with all three seeds agreeing; at $p_0 = 0.067$ it does not measurably help, under
the same harness, optimizer, budget and seeds (\S\ref{sec:p0}). The swing between them
is $0.102$, about $2.5\sigma$, and it is what carries the answer rather than either endpoint. The mechanism named the
condition in advance: unconditional refinement is nearly free where $p_0$ is small and costly
where it is not, and the evolved gene's \texttt{PASS} rate duly moves from zero to non-zero as
$p_0$ rises while never emitting a false \texttt{PASS}. \emph{The null was real and it was local}
--- every instruction-following run we bought sat at the low-$p_0$ end of an axis we had not
identified, and the corrected re-run measures it at $0.067$.

\textbf{Does GEPA lift this task?} No, at the
budgets tested: its genomes gave no held-out gain and sometimes \emph{degraded} both arms. That is
not a verdict on GEPA, which reports gains elsewhere, and it is not train-split overfitting ---
that was our reading until the sweep varied the split from $4$ to $24$ and found no trend to
overfit along (\S\ref{sec:splitresult}). Nor is it under-search: a second pre-registered sweep
spans a measured six-fold depth range with no ordering and the deepest cell flattest, while
acceptance rises from $4/6$ to $6/6$ arm-seed pairs (\S\ref{sec:budgetresult}). That dissociation
is what separates \emph{the optimizer failed to search} from \emph{the optimizer searched and
found nothing that generalizes}, and it points firmly at the second.

Four qualifications travel with the answer. At $n=3$ per level no interval is significant, since
the sign test cannot fall below $p=0.25$, so the load-bearing evidence is the registered swing
rather than any one interval. Two of our registered predictions \emph{failed}: monotonicity in
$p_0$, because the highest level saturates, and the sign of $\dd$ at $n_{\mathrm{train}}{=}24$,
which came out positive (\S\ref{sec:splitresult}). $\dd$ is a whole-system contrast: decomposed,
only $40\%$ of the $-0.042$ is the treated component while all of the $+0.060$ is
(\S\ref{sec:decomp}). And the surviving picture is narrower than the one we set out with and
better supported: the optimizer reliably
drives the verifier to never accept when $p_0$ is low, that behaviour is cheap to screen for, and
whether it helps or hurts held-out score depends on $p_0$ by about $0.1$ across the range we
bought --- against a measurement noise we can now size at $\sigma(\dd) = 0.029$
(\S\ref{sec:threats}), which every interval in the study omits.

One distinction is worth stating separately, because we had run two claims together and only one
survives. \textbf{A degenerate verifier and a harmful verifier are different objects.} A critique
that never accepts is behaviourally just \emph{always refine}, and where almost every draft is
wrong, always-refine is close to the right policy. That is why the \emph{behaviour} replicates on
$11/12$ evolved genes across four split sizes while any \emph{cost} it imposes does not replicate
at all --- three of four split-size cells are positive. We withdraw the cost claim and keep the
behavioural one, and the practical consequence is a screen rather than a prohibition: a
\texttt{PASS} rate of zero tells you the search went to the degenerate corner, and whether that
matters depends on $p_0$.

Reaching this answer cost every contrast run which preceded it: \textbf{fifteen} of them, across
instruction following, a cross-domain replication on frontier-hard code (\S\ref{sec:bcb}), an
actuator variant and a worker-model swap. Ten reported $\dd(D{-}B)$ of exactly zero, and none of
them tested the hypothesis, because the critique gene never changed:
a zero replicated ten times, from an experiment that did not run. Paper~III reports that audit and
the reporting minimum it argues for --- \textbf{per component, whether the optimizer's output
differs from the seed} --- which would have ended this line of work fifteen runs earlier.
Against that background the honest summary of what this paper contributes about componentized
optimization is one configuration in which the treatment demonstrably ran and its effect is
conditional on a measurable property of the task distribution. For fifteen runs the answer was
``nothing yet''; the sixteenth is a first data point, not a verdict.

We release the corrected harness, every cassette and pre-registration including the superseded
ones. This paper and its figures are licensed CC BY 4.0 and the accompanying code Apache-2.0, so
the artifact can be redistributed and built on rather than merely read.

\bibliographystyle{plainnat}
\bibliography{references}

\appendix

\section{Withdrawals, Corrections and Registered Deviations}
\label{sec:corrections}

This paper's earlier drafts published claims it now withdraws, and two of its pre-registered
predictions failed. Reporting guidance for preregistered work asks that deviations be disclosed in
one systematic place rather than left for a reader to reconstruct by comparing documents
\citep{willroth2024bestlaid}, so every such change affecting \emph{this} paper's claims is
collected here and the body text states only the surviving claim. \textbf{Twelve} entries follow.
Substantive withdrawals give the claim as
published and what replaces it; arithmetic corrections are grouped by cause, because the useful
information is the mechanism that produced a class of wrong digits rather than each digit, and the
superseded values themselves are recoverable from the repository history released with this paper.
The series' other two papers carry the entries bearing on their own claims.

\begin{enumerate}\itemsep2pt
  \item \textbf{The cost of evolving the verifier at low $p_0$.} Published: evolving the verifier
    alongside the generator \emph{costs} $0.042$ held-out. Withdrawn. Three further
    training-split sizes at the same $p_0$ come out positive
    (\S\ref{sec:splitresult}), the cell is $1.44\sigma$ of a measured noise floor
    (\S\ref{sec:threats}), and five independent draws of the same genomes halve it to $-0.021$,
    $1.12$ standard errors from zero (\S\ref{sec:repeat}). What survives is the
    \emph{conditional} statement in \S\ref{sec:p0}.
  \item \textbf{``Optimization made both arms worse.''} Published from the contaminated first
    attempt, where prompt-only and loop-level both fell below baseline. On the clean re-run
    arm~B is exactly flat, so the generalization does not survive; the loss is specific to the
    arm carrying the extra gene.
  \item \textbf{The outcome-side half of our own mechanism.} Published: the non-accepting
    verifier loses by damaging drafts that needed no repair. Stratified on a pre-treatment
    covariate, the already-correct stratum contributes $+0.002$ (\S\ref{sec:decomp}), so that
    half is withdrawn. The training-signal half, which predicted the sign before the runs, stands.
  \item \textbf{$\dd$ read as a treatment effect.} $\dd(D-B)$ is a whole-system contrast, because
    both arms evolve the shared generation gene. Decomposed, $60\%$ of the $-0.042$ is that
    shared gene while $100\%$ of the $+0.060$ is the treatment (\S\ref{sec:decomp}).
  \item \textbf{Registered prediction: monotonicity in $p_0$.} \emph{Falsified.} The highest
    level saturates, which was a design error rather than a property of the axis
    (\S\ref{sec:p0}).
  \item \textbf{Registered prediction: the sign of $\dd$ at $n_{\mathrm{train}} = 24$.}
    \emph{Falsified.} It came out positive (\S\ref{sec:splitresult}). Reported as a failed sign
    claim rather than retreated to its interval.
  \item \textbf{The registered search-depth formula.} It priced every proposal a reflection
    minibatch \emph{plus} a full valset re-evaluation, but a proposal rejected on the minibatch
    pays only the minibatch. Three sweep cells therefore landed on one effective depth. The
    registered text is left unedited and depth is now \emph{measured}, by replaying each cassette
    and counting the optimizer's iterations (\S\ref{sec:budgetresult}).
  \item \textbf{The $L_3$ null.} Published as $-0.002$, a single draw. Five replicates put it at
    $+0.019$ (\S\ref{sec:repeat}): still not significant, but its reported \emph{sign} was a
    sampling artifact.
  \item \textbf{Cost figures.} The harness records a pre-flight \emph{bound} on spend, not an
    invoice (\S\ref{sec:prereg}). Bounds are labelled as bounds throughout; an earlier draft
    reported one as a measured spend.
  \item \textbf{The five-replicate $t$-tests used a population standard deviation.} Every
    $K{=}5$ standard error divided the \emph{population} sd of the five draws by $\sqrt{5}$, and
    the result was then compared against critical values from $t(4)$ --- a distribution that
    exists \emph{because} the denominator is $n-1$. The unbiased estimator inflates $\sigma$ by
    $\sqrt{5/4} = 1.118$, so every $t$ in earlier drafts was $10.6\%$ too large. The same
    mistake in the cross-run estimate, over $30$ paired tasks, moved $\sigma(\dd)$ from $0.0291$
    to $0.0286$ and every ratio scaled by it by $1.7\%$. Both are corrected throughout. The
    error was invisible to \texttt{verify\_paper\_claims.py}, which checks the paper against the
    same function that computed it, so paper and checker were wrong together --- the same
    shared-error blind spot as the borrowed $\sigma$, one level down.
    It ran in the flattering direction and it moved one verdict: the $L_2$ gain was published at
    $p<0.05$ ($t = 3.02$) and is $p<0.10$ ($t = 2.70$). \S\ref{sec:repeat} states it that way, and
    the paper no longer claims any effect at $p<0.05$.
  \item \textbf{Arm~C's per-arm proposal count.} An earlier draft reported that GEPA proposed
    ``four candidates on every seed'' for the genome-blind control. Withdrawn as not derivable from
    the release: the cassette is content-addressed on a call's name and arguments, so an entry
    carries no arm label and no per-arm proposal count can be recovered from it. \S\ref{sec:rerun}
    now reports the run-level count, which is derivable and makes the same point.
  \item \textbf{Two results of the companion papers, over-claimed here.} The $+0.20$
    actionability lift on frontier-hard code was described as \emph{established}, which is stronger
    than Paper~II claims for it: that paper reports it without an interval, because a seed on
    BigCodeBench-Hard re-partitions a fixed pool rather than drawing a fresh sample. And two
    inferences drawn from runs 4 and 5 did not say that those runs predate the harness correction,
    so their between-arm contrast is not a measurement. Both are now stated at the strength the
    source supports.
\end{enumerate}

\section{The Actuator Experiment, and the MBPP+ Attempt}
\label{sec:actuator}

Two negative sub-results from the code domain, summarized in \S\ref{sec:bcb} and reported
here in full.

\paragraph{The MBPP+ attempt that motivated the move to a frontier-hard benchmark.}
The MBPP+ attempt used the base tests as
the in-loop verifier and the edge-case tests as the hidden scorer. Every arm returned the
seed genome's held-out score on both seeds, so $\dd(D{-}B)=0$ identically --- a degenerate
artifact rather than a null: GEPA logged \texttt{``All subsample scores perfect for parent
0. Skipping''} on every iteration, because the seed loop, aided by executable base-test
feedback, already solved all six training tasks. Held-out headroom looked adequate
($0.80$--$0.90$); the \emph{training} split had none, and that is what gates GEPA. We did
not weaken the worker to manufacture headroom, because prompt-optimization and
self-refinement gains scale inversely with model strength
\citep{zhou2025bbopromptscale,zhang2024smallverifiers}, so a positive $\dd(D{-}B)$ bought
that way would attenuate toward zero at frontier scale; moving headroom into the
\emph{tasks} keeps the worker in the same capability-relative regime a frontier model
occupies on frontier-hard work. The run reproduces bit-for-bit and its cassette is released
with the others.

\paragraph{Can GEPA \emph{discover} the actionable feedback? The actuator experiment.}
The $+0.20$ lift Paper~II reports indicates that the fixer's feedback \emph{matters} --- but that
feedback was hand-designed. The sharpest test of our thesis is whether GEPA can \emph{discover} the
gain by optimization when pointed squarely at it. We add a fourth evolvable gene, an
\texttt{actuator}: an LLM step that, on a failed critique, turns the visible-test signal and
the visible-test source (the same \emph{hi} material) into the repair guidance fed to the
refine step. Arm~B still evolves only \texttt{generate}; arm~D now co-evolves \texttt{generate}
\emph{and} \texttt{actuator} with the detector (\texttt{ensure\_correct}) held \emph{fixed}, so
any $D-B$ gap isolates optimization of the fixer's feedback --- the one component we know can
help. A first run ($2$ seeds, $8$/$10$) gave a flat $\dd(D-B)=0.00$: GEPA proposed sensible
actuator rewrites (e.g.\ ``produce concrete, draft-specific repair instructions; do not write
vague advice such as `fix the syntax' ''), but accepted none that beat even the training
subsample, and the single accepted \texttt{generate} rewrite did not transfer.

A \emph{powered} replication ($3$ seeds, $20$ train / $20$ held-out, $\mathrm{mmc}=50$)
reproduced the null. Prompt-only arm~B now transfers to held-out (mean $+0.067$), while arm~D
does worse: $\dd(D-B) = -0.10$ (per-seed $-0.15, -0.15, 0.00$; exact-resample $95\%$ CI
$[-0.150, +0.000]$, which \emph{includes} zero; sign $p = 0.50$). \textbf{Two corrections
apply to this paragraph as first written} (Paper~III): we reported the interval as
$[-0.150,-0.050]$, ``excluding zero'', which was a percentile-index error; and we explained
the deficit as arm~D over-fitting a larger search space, which the harness contradicts, since
under the default one-component-per-proposal selector arm~D searched the shared gene roughly
\emph{half} as hard as arm~B. What remains sayable is the ordering --- across all three code
configurations $D \le B$, never $D > B$ --- and that $-0.10$ on $20$ binary-scored tasks is
two task flips against a seed spread of four to five. Raising train size, resolution, seeds
and search depth surfaced no hidden positive for an arm whose extra gene the harness was not
equipped to optimize.

\begin{table}[h]
\centering
\small
\setlength{\tabcolsep}{5pt}
\begin{tabular}{lcp{0.20\textwidth}p{0.24\textwidth}}
\toprule
Result & per-seed & as published & corrected \\
\midrule
Actuator, powered ($n{=}3$) & $-0.15, -0.15, 0.00$ & $[-0.150,-0.050]$ ``excludes zero''
  & $\mathbf{[-0.150,+0.000]}$ \emph{includes zero} \\
Actionability, hi$-$lo ($n{=}3$) & $+0.090, +0.025, +0.075$ & $[+0.025,+0.085]$
  & $[+0.025,+0.090]$ (still excludes zero) \\
Null, run~3 ($n{=}5$) & $0,+0.055,0,0,-0.055$ & $[-0.033,+0.033]$ & unchanged \\
\bottomrule
\end{tabular}
\caption{The upper bound was too low by one atom. Only $n{=}3$ intervals are affected, and
the direction of the error flatters negative effects: it turns an interval that brackets
zero into one that appears to exclude it. Reproduced from Paper~III, where the
percentile-index defect is reported in full, because the first row is this appendix's result.}
\label{tab:ciexcl}
\end{table}

\section{The Hand-Written Verify Gene}
\label{sec:handgene}

This is the experiment summarized in \S\ref{sec:analysis}: the degenerate corner reached by hand rather than by search, run before the re-run was bought. Its two conditions are the first runs recorded under the corrected harness.

\medskip\noindent
Every mechanism above is measured on \emph{observed} critiques. It leaves open the
question that decides whether the loop-level lever is worth optimizing at all: is there a
better verify gene to be found? An optimizer here may only rewrite prose, so the ceiling
on what it could possibly discover is what prose can do, and that is cheap to bound
directly. We wrote a verify gene by hand against this paper's own findings --- it forces an
explicit measurement before each verdict, addresses false passes and false alarms
separately, and demands measured-value / target-value / smallest-edit repair text rather
than a restatement of the rule --- and ran arm~A twice on multi-hard, seed~0, $32$ held-out
tasks, changing \emph{only} the \texttt{ensure\_correct} gene. The two conditions share one
content-addressed cassette, so each task's first draft is byte-identical across them and
each pair isolates the gene exactly: same draft in, different critique, different repair.
The analysis was pre-registered and committed before the run was paid for
(\texttt{scripts/analyze\_c1.py}); total cost \$0.043.

Paired by task, the hand-written gene does not help: $n{=}32$, $17$ ties, better on $6$
tasks and worse on $9$, mean $\dd = -0.033$ with a bootstrap $95\%$ interval of
$[-0.095,+0.028]$ and sign test $p=0.61$. Our pre-registered rule returns
\textbf{inconclusive}, and deliberately not ``no effect'': the $90\%$ interval extends past
the $\pm0.05$ SESOI, so equivalence fails too. On its own that is an underpowered null.

What the same cassette shows about the \emph{detector} is sharper (Table~\ref{tab:c1}). The
hand-written gene did change its behaviour, and markedly: it emitted PASS zero times
against three, sending $32/32$ tasks to a refine round against $29/32$. But the change is
confined to one axis. Neither gene produced a single true positive on the five compliant
drafts in this split, exactly as the intrinsic critique of run~3 managed one true positive
out of five while emitting twelve false ones. Both sit at the chance diagonal; the
hand-written gene slid \emph{along} it, toward predicting negative always. What the
threshold move itself bought is small and legible: on the three tasks whose false pass it
correctly caught, the score rose $0.800 \to 0.867$ (one repaired to $1.0$, two unchanged),
worth $+0.006$ on the 32-task mean; on the five tasks whose draft was already compliant the
mean is $-0.100$, and since \emph{both} genes demanded fixes on all five, that is one damaged
draft rather than a false-alarm difference. Those eight tasks are the only ones on which the
two genes' \emph{verdicts} differ, and together they account for $-0.009$ of the $-0.033$: the
remaining $-0.023$ falls on the twenty-four tasks where both genes demanded a fix and only the
repair \emph{prose} differed. So eliminating false passes bought a real gain an order of
magnitude below the noise it sits in, and the larger part of the movement is not the threshold
at all --- it is the rewritten repair text, which is the lever Paper~II measures
directly.

\begin{table}[h]
\centering
\small
\begin{tabular}{lcccccc}
\toprule
Critique gene & $n$ & base rate & tp & fp & fn & PASS rate \\
\midrule
run 3, intrinsic (F$_1 = 0.11$) & 50 & 0.10 & 1 & 12 & 4 & 0.26 \\
C1, seed gene & 32 & 0.16 & 0 & 3 & 5 & 0.09 \\
C1, hand-written gene & 32 & 0.16 & 0 & 0 & 5 & \textbf{0.00} \\
\bottomrule
\end{tabular}
\caption{The hand-written gene changed \emph{when} the detector fires, not \emph{whether it
can tell}. F$_1$ is undefined for both C1 conditions because neither has a true positive.
With only five positive cases per split this cannot separate ``no discrimination'' from
``weak discrimination'' --- the claim is narrower: the measurable change was confined to the
threshold.}
\label{tab:c1}
\end{table}

This is a mechanism for the loop-level null that does not involve the optimizer at all. If
rewriting this gene's prose moves where its threshold sits without giving it discriminative
power, then a search over prose searches a space whose reachable improvements are threshold
moves, and a threshold move in a detector that cannot discriminate has no systematic gain to
offer. It is evidence, not proof: one hand-written point in prose space does not establish
that no prose adds discrimination, and a competent author might do better. It does mean a
loop-level null should no longer be read primarily as a fact about optimizers. This is what
made the corrected-harness re-run worth buying and what it was pre-registered to predict
(\S\ref{sec:directions}); \S\ref{sec:rerun} reports the outcome.

\end{document}
